% 地球（综述）
% license CCBYSA3
% type Wiki

本文根据 CC-BY-SA 协议转载翻译自维基百科\href{https://en.wikipedia.org/wiki/Earth}{相关文章}。

\begin{figure}[ht]
\centering
\includegraphics[width=8cm]{./figures/5280dfaa768e0550.png}
\caption{“蓝色弹珠”，阿波罗 17 号，1972 年 12 月} \label{fig_Earth_1}
\end{figure}
地球是距太阳的第三颗行星，也是唯一已知能够孕育生命的天体。这得益于地球是一颗海洋世界——在太阳系中，唯一能够维持液态地表水的星球。地球上几乎所有的水都存在于全球性海洋中，覆盖了地球地壳的 70.8\%。剩余的 29.2\% 地壳为陆地，其中大部分集中在地球陆半球的大陆地块上。地球的大多数陆地区域至少具有一定湿度并被植被覆盖，而地球极地沙漠中的大型冰盖所储存的水量超过了地球地下水、湖泊、河流和大气水的总和。地球的地壳由缓慢移动的构造板块组成，彼此相互作用形成山脉、火山与地震。地球拥有液态外核，可产生磁层，能够偏转大部分具有破坏性的太阳风和宇宙射线。

地球拥有一个动态的大气层，这维持了地球的表面环境，并在陨石体及紫外线入射时提供保护。大气层主要由氮气和氧气组成。水汽在大气层中广泛存在，形成覆盖地球大部分区域的云层。水汽是一种温室气体，并与大气中的其他温室气体（特别是二氧化碳 $CO_2$）一道，通过捕获来自太阳光的能量，创造出液态表面水和水汽能够持续存在的条件。这一过程维持了目前 14.76 °C（58.57 °F）的平均地表温度，在正常大气压下，水在该温度呈液态。不同地理区域所捕获能量的差异（如赤道地区获得的阳光多于极地地区）驱动了大气和海洋环流，形成具有不同气候区的全球气候系统，并产生降水等多种天气现象，使碳和氮等组成成分得以循环。

地球呈近似椭球形，周长约为 40,000 公里（24,900 英里）。它是太阳系中密度最高的行星。在四颗类地行星中，地球最大且质量最高。地球距太阳约八光分（1 天文单位），沿轨道绕太阳运行一周约需一年（约 365.25 天）。地球自转一周所需时间略少于一天（约 23 小时 56 分钟）。地球自转轴相对于其绕太阳轨道平面的垂线有倾斜，从而产生四季更替。地球拥有一颗永久天然卫星——月球。月球以 384,400 公里（238,855 英里，约 1.28 光秒）的距离绕地球运行，其直径约为地球的四分之一。月球的引力有助于稳定地球自转轴，产生潮汐，并逐渐减缓地球自转。同样，地球的引力已使月球自转潮汐锁定，月球永远以同一面朝向地球。

地球与太阳系中大多数其他天体一样，约在 45 亿年前由早期太阳系中的气体和尘埃形成。在地球历史的最初十亿年中，海洋形成，随后生命在海洋中出现。生命在全球扩散，并持续改变地球的大气和地表，导致约 20 亿年前发生“大氧化事件”。人类约 30 万年前在非洲出现，随后扩散到地球的每一块大陆。人类依赖地球的生物圈和自然资源而生存，但对地球环境的影响却日益增加。如今，人类对地球气候与生物圈的影响已不可持续，威胁着人类自身及许多其他生命形式的生计，并导致物种广泛灭绝。

地球呈椭球形，周长约 40,000 公里（24,900 英里）。它是太阳系中密度最高的行星。在四颗岩质行星中，地球最大且质量最高。地球距离太阳约八光分（1 个天文单位），绕太阳公转一周约需一年（约 365.25 天）。地球绕自身轴心自转一周时间略少于一天（约 23 小时 56 分钟）。地球自转轴相对于其绕太阳轨道平面垂线存在倾斜，从而产生四季变化。地球拥有一颗永久天然卫星——月球，月球以 384,400 公里（238,855 英里，即 1.28 光秒）的距离绕地运行，其直径约为地球的四分之一。月球的引力有助于稳定地球自转轴、引发潮汐，并逐渐减缓地球自转。同样，地球的引力已使月球自转产生潮汐锁定，月球始终以同一面朝向地球。

地球与太阳系中大多数其他天体一样，约在 45 亿年前由早期太阳系中的气体和尘埃形成。在地球历史最初的十亿年中，海洋开始形成，随后生命在其中诞生。生命在全球扩散，并持续改变地球的大气层和地表，导致约二十亿年前发生“大氧化事件”。人类约在 30 万年前出现在非洲，并扩散至地球的每一块大陆。人类依赖地球的生物圈和自然资源生存，但对地球环境的影响日益加剧。目前人类对地球气候和生物圈的影响已不可持续，威胁着人类和许多其他生命形式的生存，并正在引发广泛的物种灭绝。
\subsection{词源}
现代英语单词 earth 经由中古英语发展而来，源自古英语名词 eorðe $^{[22]}$。在所有日耳曼语言中都有其同源词，由此重建出原始日耳曼语形式 erþō。在最早的文献中，单词 eorðe 用来翻译拉丁语 terra 和希腊语 gē 的多重含义：地面、泥土、陆地、人类世界、世界表面（包括海洋）以及整个地球本身。与罗马神话中的 Terra（或 Tellus）和希腊神话中的 Gaia 类似，“地球”可能在日耳曼多神教中也被视为一位拟人化的女神：晚期北欧神话中出现过 Jörð（“大地”），这位女巨人常被认为是雷神托尔的母亲 $^{[23]}$。

历史上，“earth”通常以小写书写。在中古英语早期，为表达“地球”这一明确含义，开始使用短语 the earth。到了早期现代英语时期，名词大写逐渐普及，the earth 也写作 the Earth，尤其是在与其他天体并列提及时。近来，这一名称有时直接写作 Earth，以与其它行星名称类比，然而 earth 和 the earth 的用法仍然常见 $^{[22]}$。不同出版风格各有规定：牛津拼写认为小写形式最为常见，而大写形式则为可接受的变体。另一种惯例是在作为名称使用时采用大写（如 “Earth's atmosphere”），但在前有定冠词时使用小写（如 “the atmosphere of the earth”）。在口语表达中几乎总是写作小写，例如“what on earth are you doing?” $^{[24]}$

名称 Terra（/ˈtɛrə/）有时用于科学写作；它也常见于科幻作品中，用来将人类居住的行星与其他星球区分开来 $^{[25]}$。而在诗歌中，Tellus（/ˈtɛləs/）则被用来表示地球的拟人形象 $^{[26]}$。Terra 也是一些拉丁语系语言（如意大利语和葡萄牙语）中“地球”的名称；在其他罗曼语中，该词的拼写略有变化，如西班牙语 Tierra 与法语 Terre。希腊诗意名称 Gaia 的拉丁化形式 Gaea 在英语中较为少见，尽管由于“盖亚假说（Gaia hypothesis）”，另一种拼写 Gaia 已变得常用，并读作 /ˈɡaɪ.ə/，而不是更传统的英语发音 /ˈɡeɪ.ə/ $^{[27]}$。

关于地球还有许多形容词。earthly（地上的、尘世的）源自 Earth。由拉丁语 Terra 衍生出 terran（/ˈtɛrən/）、terrestrial（/təˈrɛstriəl/），以及（经由法语） terrene（/təˈriːn/）；由拉丁语 Tellus 衍生出 tellurian（/tɛˈlʊəriən/）和 telluric $^{[28][29][30][31][32]}$。
\subsection{自然历史}
\subsubsection{形成}
\begin{figure}[ht]
\centering
\includegraphics[width=8cm]{./figures/409cc725593ef0b7.png}
\caption{对形成地球及其他太阳系天体的早期太阳系原行星盘的描绘} \label{fig_Earth_2}
\end{figure}
在太阳系中发现的最古老物质被测定为距今 $4.5682^{+0.0002}_{-0.0004}
$ 十亿年前$^{[33]}$。在距今 $4.54\pm0.04$ 十亿年前，原始地球已经形成$^{[34]}$。太阳系中的各个天体与太阳一起形成并演化。理论上，太阳星云通过引力坍缩从分子云中分割出一块空间，开始旋转并扁平化成一个环星盘，然后行星在该盘内与太阳一起生长。星云包含气体、冰粒和尘埃（包括原始核素）。根据星云理论，微行星通过吸积而形成，据估计原始地球可能需要约 7000 万年至 1 亿年才能形成$^{[35]}$。

对月球年龄的估计从距今 45 亿年到显著更年轻不等$^{[36]}$。主流假说认为，月球由一次撞击后从地球上抛出的物质吸积而成——一个约为地球质量 10\%、大小与火星相似的天体（命名为 忒伊亚 Theia）与地球发生碰撞$^{[37]}$。它以掠擦方式撞击地球，其部分质量与地球合并$^{[38][39]}$。大约在距今 40 亿年至 38 亿年前，晚期重轰炸**期间大量小行星撞击导致月球大范围表面环境发生显著改变，并据此推断地球当时的表面环境也发生了重大变化$^{[40]}$。
\subsubsection{形成之后}
地球的大气和海洋由火山活动和气体逸出形成$^{[41]}$。来自这些源的水蒸气凝结成海洋，并由小行星、原行星和彗星带来的水与冰进一步补充$^{[42]}$。按照一种模型，自地球形成之初起，地球上就可能已经有足以填满海洋的水量$^{[43]}$。在该模型中，大气中的温室气体阻止了当时只有当前 70\%光度的新生太阳导致海洋结冰$^{[44]}$。到距今约 35 亿年前，地球的磁场已经建立，这有助于防止大气被太阳风剥离$^{[45]}$。
\begin{figure}[ht]
\centering
\includegraphics[width=8cm]{./figures/e931dfc771eed8c4.png}
\caption{“淡橙色小点——对早期地球的一种印象画，呈现其富含甲烷并带有橙色色调的早期大气层$^{[46]}$。”} \label{fig_Earth_3}
\end{figure}
随着地球熔融的外层冷却，形成了第一批固体地壳，被认为成分为镁铁质$^{[47]}$。第一批大陆地壳成分更为长英质，是由这种镁铁质地壳部分熔融形成的$^{[47]}$。在太古宙沉积岩中发现的冥古宙时代锆石颗粒表明，至少部分长英质地壳早在 44 亿年前就已经存在，比地球形成仅晚约 1.4 亿年$^{[48]}$。

关于这一初始、体积很小的大陆地壳如何演化至如今的丰富状态，主要有两种模型$^{[49]}$：
(1) 相对稳定地持续增长直至今天$^{[50]}$，这一观点得到全球大陆地壳放射性定年数据的支持；
(2) 在太古宙期间大陆地壳体积迅速增长，形成了现今大陆地壳的主体$^{[51][52]}$，这一观点得到来自锆石中铪同位素和沉积岩中钕同位素证据的支持。通过大陆地壳的大规模循环再造，特别是在地球早期阶段，这两种模型及其证据是可以调和的$^{[53]}$。

新的大陆地壳是通过板块构造形成的，而板块构造最终由地球内部持续散失热量驱动。在数亿年的时间尺度上，构造力促使大陆地壳区域结合成超级大陆，随后又破裂分离。约在 7.5 亿年前，最早已知的超级大陆之一罗迪尼亚（Rodinia）开始裂解。之后大陆重新聚合，在 6 亿至 5.4 亿年前形成泛诺西亚（Pannotia），最终又形成盘古大陆（Pangaea），该超级大陆在约 1.8 亿年前开始裂解$^{[54]}$。

最近一次冰期循环模式始于约 4,000 万年前$^{[55]}$，并在约 300 万年前的更新世显著增强$^{[56]}$。自那以后，高纬度和中纬度地区经历了反复的冰川与间冰期循环，大约每 21,000 年、41,000 年和 100,000 年重复一次$^{[57]}$。最近一次冰期，俗称“最后冰期（Last Glacial Period）”，使大陆大部分区域至中纬度地区都被冰层覆盖，并于约 11,700 年前结束$^{[58]}$。
\subsubsection{生命起源与演化}
化学反应在约 40 亿年前产生了第一批能够自我复制的分子。约 5 亿年后，所有现存生命的最后共同祖先出现$^{[59]}$。光合作用的进化使生命形式能够直接利用太阳能。随之产生的分子氧（O(_2)）在大气中累积，并在与太阳紫外辐射的相互作用下，于高层大气中形成了保护性的臭氧层（O(_3)）$^{[60]}$。较小细胞被纳入较大细胞之中，产生了复杂的细胞——真核细胞$^{[61]}$。随着群体内细胞分化程度不断提高，真正的多细胞生物得以形成。在臭氧层吸收有害紫外辐射的帮助下，生命开始殖民地球表面$^{[62]}$。最早的生命化石证据包括：在西澳大利亚约 34.8 亿年前的砂岩中发现的微生物席化石$^{[63]}$；在格陵兰西部约 37 亿年前的变沉积岩中发现的生物成因石墨$^{[64]}$；在西澳大利亚约 41 亿年前岩石中发现的生物材料残留物$^{[65][66]}$。地球上最早的直接生命证据存在于约 34.5 亿年前的澳大利亚岩石中，其中显示了微生物的化石$^{[67][68]}$。
\begin{figure}[ht]
\centering
\includegraphics[width=8cm]{./figures/5a5b88da1cf2e867.png}
\caption{对太古宙（Archean）的艺术印象——这是地球形成之后的那个宙——画面中出现了圆形叠层石，它们是数十亿年前产生氧气的早期生命形式。晚期重轰炸（Late Heavy Bombardment）结束后，地壳已经冷却，富含水分的贫瘠地表上可以看到大陆和火山的痕迹，而当时的月球仍以今天一半的距离绕地球运行，看起来大约大 2.8 倍，并产生强烈的潮汐$^{[69]}$。} \label{fig_Earth_4}
\end{figure}
在新元古代（Neoproterozoic），约公元前 10 亿年至 5.39 亿年前（1000–539 Ma），地球表面可能有相当大部分被冰层覆盖。这个假说被称为“雪球地球”（Snowball Earth），它之所以备受关注，是因为这一时期发生在寒武纪大爆发之前，而寒武纪大爆发（约 5.35 亿年前，535 Ma）标志着多细胞生命复杂性显著增加的时期$^{[70][71]}$。

在寒武纪大爆发之后，地球上至少发生过五次大型物种大灭绝事件，以及众多较小的灭绝事件$^{[72]}$。除目前被提出的全新世灭绝（Holocene extinction event）之外，最近的一次发生在约 6600 万年前（66 Ma），当时一颗小行星撞击引发了非鸟类恐龙及其他大型爬行动物的灭绝，但小型动物——例如昆虫、哺乳动物、蜥蜴和鸟类——在很大程度上得以幸存。

在过去的 6600 万年（66 Mys）中，哺乳动物不断多样化；数百万年前，非洲的一种类人猿获得了直立行走的能力$^{[73][74]}$。直立行走使得工具使用成为可能，并促进交流，从而为更大的大脑提供所需的营养与刺激，进而推动了人类的演化。

农业的发展，随后是文明的出现，使人类开始深刻影响地球以及其他生命形式的性质与数量——这种影响一直持续至今$^{[75]}$。
\subsubsection{未来}
\begin{figure}[ht]
\centering
\includegraphics[width=8cm]{./figures/a4705ee315f9141e.png}
\caption{对大约 50–70 亿年后，当太阳进入红巨星阶段时灼热地球的构想图} \label{fig_Earth_5}
\end{figure}
地球预期的长期未来与太阳的未来紧密相关。在未来约 11 亿年内，太阳光度将增加 10\%，在未来约 35 亿年内将增加 40\%$^{[76]}$。地球表面温度的上升将加速无机碳循环，可能在约 1 亿至 9 亿年内将大气 CO(_2) 浓度降低至对现存植物致命的水平（对于 C4 光合作用约为 10 ppm）$^{[77][78]}$。植被的缺失将导致大气中氧气的消失，使现有动物生命无法存活$^{[79]}$。由于光度升高，地球的平均温度可能在 15 亿年后达到 100 ℃（212 ℉），所有海洋水分将会蒸发并逃逸到太空，从而可能在大约 16 亿至 30 亿年内触发失控温室效应$^{[80]}$。即便太阳保持稳定且永恒存在，由于随着地核缓慢冷却，中洋脊蒸汽喷发减少，现代海洋中的相当一部分水也将会沉入地幔$^{[80][81]}$。

太阳将在约 50 亿年后演化成红巨星。模型预测太阳将膨胀至大约 1 AU（1.5 亿千米；9,300 万英里），约为其当前半径的 250 倍$^{[76][82]}$。地球的命运尚不清楚。作为红巨星，太阳将失去大约 30\% 的质量，因此，如果没有潮汐效应，当太阳达到最大半径时，地球将移动到距离太阳 1.7 AU（2.5 亿千米；1.6 亿英里）的轨道；否则，在潮汐效应存在的情况下，它可能进入太阳大气并被汽化，其较重的元素沉入垂死太阳的核心$^{[76]}$。
\subsection{物理特性}
\subsubsection{大小与形状}
\begin{figure}[ht]
\centering
\includegraphics[width=8cm]{./figures/a3439116fb51e49f.png}
\caption{地球西半球，显示的是相对于地心的地形高度，而不是像常见地形图那样以平均海平面为参考。} \label{fig_Earth_6}
\end{figure}
地球由于处于流体静力平衡状态，呈现圆润的形状，平均直径为 12,742 公里（7,918 英里），是太阳系中第五大类地行星尺度天体，也是最大的类地天体。$^{[83][84]}$

由于地球自转，它呈现扁球体形状，在赤道发生隆起；其赤道直径比两极直径长 43 公里（27 英里）。$^{[85][86]}$ 地球的形状还存在局部地形变化；最大的局部变化，如马里亚纳海沟（低于当地海平面 10,925 米，约 35,843 英尺）会使地球平均半径缩短约 0.17\%，而珠穆朗玛峰（高于当地海平面 8,848 米，约 29,029 英尺）会使地球平均半径增加约 0.14\%。$^{[\text{n}6][89]}$ 由于地表距离地球质心在赤道隆起处最远，因此厄瓜多尔的钦博拉索火山顶（6,384.4 公里，或 3,967.1 英里）是地球表面距离地心最远的点。$^{[90][91]}$ 与较为刚性的陆地地形类似，海洋也呈现更具动态的地形变化。$^{[92]}$

为了测量地球局部地形的变化，大地测量学采用一种理想化的地球模型，得到大地水准面（geoid）形状。这种形状是通过将海洋理想化处理，即假设其完全覆盖地球且没有潮汐、风等扰动而得到的。其结果是一个光滑但不规则的大地水准面，提供一个平均海平面作为地形测量的参考高度。$^{[93]}$
\subsubsection{表面}
\begin{figure}[ht]
\centering
\includegraphics[width=8cm]{./figures/958849ebee4464c0.png}
\caption{地球的合成影像，清晰可分辨其不同类型的表面：海洋（蓝色）占据地球表面的主导区域；非洲大陆从葱郁（绿色）到干旱（棕色）的土地分布明显；而地球极地冰层以南极海冰（灰色）覆盖南大洋，以及南极冰盖（白色）覆盖南极大陆。} \label{fig_Earth_7}
\end{figure}
地球表面是大气层与固体地壳及海洋之间的边界。按此定义，其面积约为 5.1 亿平方千米（1.97 亿平方英里）$^{[11]}$。地球可按纬度分为北半球和南半球，也可按经度分为东半球和西半球。

地球表面的大部分为海洋水体：70.8\%，即 3.61 亿平方千米（1.39 亿平方英里）$^{[11]}$。这片巨大的咸水体通常称为世界海洋$^{[94][95]}$，使得拥有动态水圈的地球成为一个真正的水之星球$^{[96][97]}$，或海洋世界$^{[98][99]}$。事实上，在地球早期历史中，海洋可能曾完全覆盖地球$^{[100]}$。世界海洋通常被划分为太平洋、大西洋、印度洋、南冰洋和北冰洋从最大到最小依次排列。海洋覆盖着地球的大洋地壳，其中大陆架海在较小程度上覆盖着大陆地壳的大陆架。大洋地壳形成了大型的海洋盆地，具有诸如：深海平原,海底山,海底火山 $^{[85]}$,海沟,海底峡谷,大洋高原,以及贯穿全球的大洋中脊系统$^{[101]}$在地球极地区域，海洋表面覆盖着季节性变化的海冰，这些海冰常与极地陆地、永久冻土及冰盖相连，形成极地冰盖。
\begin{figure}[ht]
\centering
\includegraphics[width=8cm]{./figures/55e03c2e7e1d60d7.png}
\caption{地壳地形起伏} \label{fig_Earth_9}
\end{figure}
地球的陆地覆盖了地表的 29.2\%，约 $1.49\times10^{8},\text{km}^2$（5,800 万平方英里）。这些陆地包括全球范围内大量岛屿，但大多数陆地面积由四大洲块构成（按面积从大到小排列）：非洲－欧亚大陆、美洲、南极洲和澳大利亚 $^{[102][103][104]}$。这些大陆块进一步被划分并归类为各大洲。陆地表面的地形变化巨大，包括山区、沙漠、平原、高原以及其他地貌形态。陆地的海拔从最低的 −418 m（−1,371 英尺，位于死海）到最高的 8,848 m（29,029 英尺，珠穆朗玛峰峰顶）。陆地平均海拔约为 797 m（2,615 英尺）$^{[105]}$。

陆地表面可能被地表水、积雪、冰层、人造结构或植被覆盖。地球上大多数陆地具有植被覆盖 $^{[106]}$，但也有相当数量的陆地为冰盖（10\%，不包括同样面积巨大的永久冻土下的陆地 $^{[107][108]}$）或沙漠（33\% $^{[109]}$）。

土壤圈（pedosphere）是地球陆地表面最外层，由土壤构成，并受土壤形成过程影响。土壤对于陆地可耕性至关重要。地球可耕地约占陆地表面的 10.7\%，其中 1.3\% 为永久农田 $^{[110][111]}$。全球耕地面积估计约为 $1.67\times10^{7},\text{km}^2$（640 万平方英里），牧场面积约为 $3.35\times10^{7},\text{km}^2$（1,290 万平方英里）$^{[112]}$。

陆地表面与海底构成了地壳的顶部，并与上地幔部分共同形成岩石圈。地壳可分为洋壳和陆壳。洋壳在海底沉积物之下主要由玄武岩组成，而陆壳可能包含密度更低的岩石，如花岗岩、沉积岩和变质岩 $^{[113]}$。大陆表面约有 75\% 被沉积岩覆盖，尽管沉积岩仅占整个地壳质量的约 5\% $^{[114]}$。

地表地形包括海洋表面的地形及陆地形态。海底地形平均水深约 4 km，其地貌变化程度与海平面以上的陆地同样丰富。地球表面持续受到内部板块构造过程的塑造，包括地震与火山活动；以及由冰、水、风和温度驱动的风化与侵蚀；还包括生物过程，如生物量的生长与分解形成土壤 $^{[115][116]}$。
\subsubsection{板块构造}
\begin{figure}[ht]
\centering
\includegraphics[width=8cm]{./figures/c9553aaef0836668.png}
\caption{} \label{fig_Earth_10}
\end{figure}
地球地壳和上地幔构成的机械刚性外层，即岩石圈，被划分为若干构造板块。这些板块是彼此相对运动的刚性单元，并在三种类型的板块边界处发生移动：在聚合边界，两块板块相互靠拢；在张裂边界，两块板块被拉开；在转换边界，两块板块沿水平方向相互错动。在这些板块边界处，可能发生地震、火山活动、造山运动以及海沟的形成 $^{[118]}$。这些构造板块“漂浮”在软流圈之上，软流圈是上地幔中固态但黏度较低的部分，可以随板块一起流动和移动 $^{[119]}$。

随着构造板块的移动，洋壳在聚合边界处被俯冲到板块的前缘之下。同时，地幔物质在张裂边界处上涌，形成大洋中脊。这些过程共同作用，使洋壳不断回收进入地幔。由于这种循环，绝大部分海洋底部的年龄少于 1 亿年。年龄最老的洋壳位于西太平洋，估计约有 2 亿年 $^{[120][121]}$。相比之下，最早测定的陆壳年龄为 40.30 亿年 $^{[122]}$，虽然在始太古代沉积岩中保存的碎屑锆石显示年龄可达 44 亿年，表明当时至少已经存在少量大陆地壳 $^{[48]}$。

七大主要板块包括太平洋板块、北美板块、欧亚板块、非洲板块、南极板块、印－澳板块和南美板块。其他著名板块包括阿拉伯板块、加勒比板块、位于南美洲西海岸外侧的纳斯卡板块，以及南大西洋的斯科舍板块。澳大利亚板块在距今 5,000–5,500 万年前与印度板块融合。移动速度最快的是洋壳板块，例如科科斯板块以每年 75 mm（3.0 英寸/年）的速度前进 $^{[123]}$，太平洋板块以每年 52–69 mm（2.0–2.7 英寸/年）移动。相反，移动速度最慢的是南美板块，典型速度约为每年 10.6 mm（0.42 英寸/年）$^{[124]}$。
\subsubsection{内部结构}
\begin{figure}[ht]
\centering
\includegraphics[width=6cm]{./figures/ff301f5286e7d7b9.png}
\caption{} \label{fig_Earth_11}
\end{figure}
地球的内部结构与其他类地行星类似，根据化学或物理（流变学）性质被划分为不同层次。最外层是成分上独特的硅酸盐固态地壳，其下方是高黏度的固态地幔。地壳与地幔之间由莫霍不连续面（Moho面）分隔 $^{[127]}$。地壳厚度在海洋之下约为 6 公里（3.7 英里），而在大陆地区约为 30–50 公里（19–31 英里）。地壳与上地幔寒冷、坚硬的顶部共同构成岩石圈，岩石圈被分割为彼此独立运动的构造板块 $^{[128]}$。

在岩石圈之下是软流圈，它是相对低黏度的一层，岩石圈“漂浮”其上。在距地表约 410 公里和 660 公里处（250 与 410 英里），地幔内部会发生重要的晶体结构变化，形成了分隔上地幔与下地幔的过渡带。地幔以下是一层极低黏度的液态外核，其下是固态内核 $^{[129]}$。地球内核可能以略高于地球其余部分的角速度旋转，每年大约提前 0.1–0.5°，尽管也有提出更高或更低速率的模型 $^{[130]}$。内核半径约为地球半径的五分之一。密度随深度增加。在太阳系所有具行星尺度的天体中，地球是密度最高的对象。
\subsubsection{化学成分}
地球质量约为 $5.97\times 10^{24}kg$（5.970 Yg）。其主要成分为铁（按质量占 32.1\%）、氧（30.1\%）、硅（15.1\%）、镁（13.9\%）、硫（2.9\%）、镍（1.8\%）、钙（1.5\%）和铝（1.4\%），其余 1.2\% 为其他元素的微量成分。由于重力分异作用，地核主要由密度较高的元素组成：铁（88.8\%），并含少量镍（5.8\%）、硫（4.5\%）及不足 1\% 的微量元素 $^{[131][47]}$。

地壳中最常见的岩石成分是氧化物。超过 99\% 的地壳由十一种元素的各种氧化物组成，主要包括含硅的氧化物（硅酸盐矿物）以及含铝、铁、钙、镁、钾或钠的氧化物 $^{[132][131]}$。
\subsubsection{内部热量}
\begin{figure}[ht]
\centering
\includegraphics[width=8cm]{./figures/2d1016cbebea8e6a.png}
\caption{地球内部向地壳表面传热分布图，大部分热流集中在大洋中脊附近} \label{fig_Earth_12}
\end{figure}
地球内部热量的主要来源包括原始热（地球形成时遗留下来的热量）和放射性热（由放射性衰变产生的热量）$^{[133]}$。地球内部主要的产热同位素为钾-40、铀-238和钍-232$^{[134]}$。在地球中心，温度可能高达 6 000 °C（10 830 °F）$^{[135]}$，压力可达 360 GPa（5 200 万 psi）$^{[136]}$。由于大量热量来自放射性衰变，科学家推测，在地球早期历史中，在半衰期较短的同位素耗尽之前，地球的产热量要高得多。在约 30 亿年前，地球产生的热量可能是今天的两倍，这会加速地幔对流和板块构造，并促成如柯马提岩（komatiite）等如今罕见的火成岩的形成$^{[137][138]}$。

地球平均热流为 87 mW/m$^{2}$，全球总热流约为 $4.42\times 10^{13}$ W$^{[139]}$。部分来自地核的热能通过地幔柱输送至地壳，这是一种由高温岩石上涌构成的对流形式。地幔柱可以产生热点和洪流玄武岩$^{[140]}$。更多的热量则通过板块构造流失，即与海岭相关的地幔上涌过程。最后一种主要的失热方式是通过岩石圈传导，其中大部分发生在海洋下方$^{[141]}$。
\subsubsection{重力场}
参见主条目：Earth 的重力
地球重力是由于地球内部质量分布而赋予物体的加速度。在地球表面附近，重力加速度约为 $9.8 m/s^{2}$（$32 ft/s^{2}$）。地形、地质以及更深层构造的局部差异会导致地球重力场的局部及区域性变化，即所谓的重力异常$^{[142]}$。
\subsubsection{磁场}
\begin{figure}[ht]
\centering
\includegraphics[width=8cm]{./figures/4946ce6297772698.png}
\caption{地球磁层示意图，显示太阳风从左向右流动。} \label{fig_Earth_13}
\end{figure}
地球磁场的主要部分是在地核中产生的，那里发生着一种动力发电机过程，将由热和成分驱动的对流的动能转化为电能与磁场能。磁场从地核向外延伸，穿过地幔，一直延伸到地球表面，在那里大致呈偶极子形态。偶极子的两极位于地理南北极附近。在磁赤道，地表的磁场强度为$3.05\times10^{-5},\mathrm{T}$，在 2000 纪元时具有$7.79\times10^{22},\mathrm{A,m^2}$的磁偶极矩，并且每世纪大约减弱近 6\%（尽管仍然强于长期平均值）$^{[143]}$。地核中的对流运动是混沌的；磁极会漂移并周期性改变方向。这导致主磁场的长期变化，以及每隔数百万年平均发生数次的不规则磁场反转。最近一次反转发生在大约 70 万年前$^{[144][145]}$。

地球磁场在空间中的延伸范围定义了磁层。太阳风中的离子和电子会被磁层偏转；太阳风压力会压缩朝向太阳一侧的磁层至约 10 个地球半径，并把夜侧磁层延伸成一条长长的尾部$^{[146]}$。由于太阳风速度大于扰动在太阳风中传播的波速，在日侧磁层前方、太阳风内部会形成超音速弓形激波$^{[147]}$。

带电粒子被约束在磁层中；等离子层（plasmasphere）由低能粒子组成，它们基本上沿着磁力线随地球自转而运动$^{[148][149]}$；环电流（ring current）由相对于地磁场漂移的中能粒子组成，但其轨迹仍主要受磁场主导$^{[150]}$；范艾伦辐射带（Van Allen radiation belts）由高能粒子组成，其运动几乎是随机的，但仍被限制在磁层内部$^{[151][152]}$。在磁暴和次暴期间，带电粒子可能从外磁层，特别是磁尾区域偏转，沿磁力线被输送到地球电离层，在那里大气原子会被激发和电离，从而产生极光$^{[153]}$。
\subsection{轨道与自转}
\subsubsection{自转}
地球相对于太阳的自转周期——即其平均太阳日——为平均太阳时的 86,400 秒（86,400.0025 SI 秒）$^{[154]}$。由于潮汐减速作用，地球的太阳日现在略长于 19 世纪，因此每天的长度比平均太阳日长 0 到 2 毫秒不等$^{[155][156]}$。

地球相对于恒星（即固定恒星）的自转周期，由国际地球自转与参考系统服务（IERS）称为恒星日，为平均太阳时（UT1）的 86,164.0989 秒，即 $23^h56^m 4.0989^s$,$^{[2][n,10]}$。地球相对于岁差或移动的平均春分点（当太阳位于赤道 90° 时）的自转周期为平均太阳时（UT1）的 86,164.0905 秒（$23^h56^m 4.0905^s$）$^{[2]}$。因此，恒星日比岁差日短约 8.4 ms$^{[157]}$。

除了大气层内的流星和低轨卫星以外，地球天空中天体的主要视运动方向为向西，视运动速率为 15°/h = 15′/min。对于靠近天球赤道的天体而言，这相当于每两分钟一个太阳或月亮的视直径；从地球表面观察，太阳和月亮的视大小大致相同$^{[158][159]}$。
\subsubsection{轨道}
\begin{figure}[ht]
\centering
\includegraphics[width=8cm]{./figures/5ba001e8ae217704.png}
\caption{夸张示意图显示了地球绕太阳的椭圆轨道，并标明轨道极值点（远日点与近日点）与四个季节的极值点——春分、秋分与夏至、冬至——并不相同。} \label{fig_Earth_14}
\end{figure}
地球围绕太阳公转，使地球成为距离太阳第三近的行星，并属于太阳系内侧行星。地球的平均轨道距离约为 1.5 亿千米（9,300 万英里），这是天文单位（AU）的基准，约等于 8.3 分钟的光行程，或约为地月距离的 380 倍。地球每 365.2564 个平均太阳日绕太阳运行一周，即一个恒星年。由于太阳在地球天空中以每天约 $1^\circ$ 向东的视运动——相当于每 12 小时一个太阳或月亮的视直径——平均而言，地球需要 24 小时（一个太阳日）才能完成自转，使太阳重新返回子午线。

地球的平均轨道速度为 29.7827 千米/秒（107,218 千米/小时；66,622 英里/小时），这一速度足以在 7 分钟内走完地球直径约 12,742 千米（7,918 英里），并在约 3.5 小时内走完地月距离 384,400 千米（238,855 英里）$^{[3]}$。

月球与地球每 27.32 天围绕共同的质心相对于背景恒星运行一周。当与地月系统绕太阳的共同轨道相结合，从朔月到朔月的会合月周期为 29.53 天。从天球北极观测，地球、月球及其自转均呈逆时针方向运动。从太阳和地球北极上方的观测点看，地球绕太阳亦呈逆时针方向公转。轨道平面与自转轴面并非完全对齐：地球自转轴相对于地日平面（黄道面）的垂直方向倾斜约 $23.44^\circ$，而地月轨道平面相对于地日平面倾斜最多可达 $\pm5.1^\circ$。若无此倾角，每隔两周就会发生一次食现象，在月食与日食之间交替$^{[3][160]}$。

地球的希尔球（Hill sphere），亦即地球的引力影响范围半径，约为 150 万千米（93 万英里）$^{[161][n,11]}$。这是地球引力相对于更远处太阳和行星引力仍为主导的最大距离。天体必须在此半径内绕地球运行，否则可能因太阳的引力摄动而脱离轨道$^{[161]}$。地球连同整个太阳系位于银河系之中，并以距银河系中心约 28,000 光年的距离绕其运行。地球位于猎户臂之中，约在银河盘面上方 20 光年的位置$^{[162]}$。
\subsubsection{自转轴倾角与季节}
\begin{figure}[ht]
\centering
\includegraphics[width=8cm]{./figures/bd222286f4ad7ce0.png}
\caption{地球自转轴的倾斜导致地球在绕太阳轨道的不同位置上，获得不同角度的季节性照射。} \label{fig_Earth_15}
\end{figure}
地球的自转轴倾角约为 23.439281°$^{[2]}$，且按照定义，地球轨道平面的轴始终指向天极。由于地球自转轴的倾斜，抵达地球表面任意给定地点的太阳光照随一年中的时间而变化。这造成了气候的季节变化：当北回归线面向太阳时，北半球出现夏季；当南回归线面向太阳时，南半球出现夏季。在每一种情况下，另一半球同时进入冬季。

在夏季，白昼时间更长，太阳在天空中的高度更高。到了冬季，气候变得更为寒冷，白昼变短$^{[163]}$。在北极圈以北和南极圈以南的地区，每年有一段时间完全没有日照，形成极夜，并在极点本身持续数月。同样的纬度也会经历“午夜太阳”现象，即太阳整天可见$^{[164][165]}$。

根据天文学惯例，四季可通过至点——轨道中地球自转轴指向或背离太阳的极值点——以及分点——地球自转轴与公转轨道轴对齐的时刻——来确定。在北半球，冬至目前出现在大约 12 月 21 日；夏至在 6 月 21 日附近，春分约为 3 月 20 日，秋分在 9 月 22 日或 23 日左右。在南半球，情况正好相反，夏至与冬至互换，春分与秋分日期也相反$^{[166]}$。

地球自转轴倾角在长时间尺度上相对稳定。其轴倾会发生章动，即一种主要周期为 18.6 年的微小、不规则运动$^{[167]}$。地球自转轴的方向（而非倾角）也会随时间变化，每 25,800 年周期绕出一个完整的圆，这种岁差是恒星年与回归年不同的原因。这两种运动均由太阳和月球对地球赤道隆起部分的引力变化所致。地球极点在地表上也会移动数米。这种极移包含多个周期性成分，合称为准周期运动。除了这种运动的年度成分外，还有一个称为钱德勒摆动的 14 个月周期。地球自转速度也会变化，这一现象称为日长变化$^{[168]}$。

地球的年度轨道为椭圆而非圆形，其最接近太阳的位置称为近日点。在现代，地球的近日点大约发生在 1 月 3 日，而远日点大约发生在 7 月 4 日。由于岁差和轨道变化，这些日期会随时间变化，后者遵循被称为米兰科维奇周期的循环模式。地日距离在一年中的变化使得地球在近日点接收到的太阳能量比远日点增加约 6.\%$^{[169][n,12]}$。由于南半球大致在地球最接近太阳的同时倾向太阳，因此一年中南半球接收的太阳能量略多于北半球。这一效应远小于由自转轴倾角导致的总能量变化，而且大部分多余能量被南半球较高比例的水体吸收$^{[170]}$。
\subsection{地月系统}
\subsubsection{月球}
\begin{figure}[ht]
\centering
\includegraphics[width=8cm]{./figures/06a0a44355a583db.png}
\caption{在火星勘测轨道飞行器视角下看到的地球与月球} \label{fig_Earth_16}
\end{figure}
月球是一颗相对较大的类地、类行星天然卫星，其直径约为地球的四分之一。就与其行星的相对大小而言，它是太阳系中最大的卫星，尽管卡戎相对于矮行星冥王星而言更大$^{[171][172]}$。其他行星的天然卫星也因地球的月球而被称作“月亮”$^{[173]}$。关于月球起源最广为接受的理论——巨撞击假说——认为月球形成于一颗名为忒伊亚（Theia）、大小类似火星的原行星与早期地球的撞击。该假说解释了月球铁和挥发性元素的相对缺乏，以及其成分与地球地壳几乎相同这一事实$^{[38]}$。计算机模拟显示，这颗原行星的两个块状残余物可能位于地球内部$^{[174][175]}$。

地球与月球之间的引力作用在地球上引起了海潮$^{[176]}$。这一作用在月球上导致其潮汐锁定：月球的自转周期与其绕地球公转所需时间相同。因此，它总是以同一面对着地球$^{[177]}$。随着月球绕地球运行，其表面不同部分受到太阳照射，形成月相变化$^{[178]}$。由于潮汐相互作用，月球正以大约 38 mm/a（每年约 1+1/2 英寸）的速度远离地球。数百万年的微小变化——以及地球日长每年增加约 23 μs——积累起来会造成显著改变$^{[179]}$。例如，在埃迪卡拉纪（约 6.2 亿年前），一年约有 400±7 天，每天约为 21.9±0.4 小时$^{[180]}$。

月球可能通过调节地球气候，对生命的发展产生剧烈影响。古生物学证据和计算机模拟表明，地球自转轴由于与月球的潮汐相互作用而得以稳定$^{[181]}$。一些理论家认为，若无这种抵御太阳和行星对地球赤道隆起产生力矩的稳定作用，地球自转轴可能会出现混沌不稳定，在数百万年尺度上发生巨大变化，类似火星的情况，尽管这一观点存在争议$^{[182][183]}$。

从地球观测，月球的距离恰好使其视直径几乎与太阳相同。这两个天体的角大小（或立体角）之所以匹配，是因为尽管太阳直径约为月球的 400 倍，但它距离地球也约为月球的 400 倍$^{[159]}$。这使得日全食和日环食在地球上得以发生$^{[184]}$。
\subsubsection{小行星与人造卫星}
\begin{figure}[ht]
\centering
\includegraphics[width=8cm]{./figures/ae4ff0c958a95f7a.png}
\caption{一幅计算机生成的图像，展示了地球同步轨道与近地轨道中人造卫星及太空碎片分布状况的映射。} \label{fig_Earth_17}
\end{figure}
地球的同轨小行星族群由准卫星、马蹄形轨道天体以及特洛伊小行星组成。已知至少有七颗准卫星，包括 469219 Kamoʻoalewa，其直径范围从 10 m 到 5000 m 不等$^{[185][186]}$。一颗特洛伊小行星伴星 2010 TK7 正在地球绕太阳轨道的领先拉格朗日三角点 L4 附近作平动$^{[187]}$。微小的近地小行星 2006 RH120 大约每二十年会接近地月系统一次，在这些接近期间，它可能在短时间内绕地球运行$^{[188]}$。

截至 2021 年 9 月，共有 4,550 颗在轨运行的人造卫星绕地球飞行$^{[189]}$。此外还有已失效的卫星，包括目前仍在轨的最古老卫星“先锋 1 号”（Vanguard 1），以及超过 16,000 块被跟踪的太空碎片$^{[n,13]}$。地球最大的人工卫星是国际空间站（ISS）$^{[190]}$。
\subsection{水圈}
\begin{figure}[ht]
\centering
\includegraphics[width=8cm]{./figures/76246827a4922672.png}
\caption{从整体视角观察地球，其全球性海洋与云层覆盖主导了地球表面与水圈；在地球的两极地区，水圈形成了更为广阔的冰层覆盖。} \label{fig_Earth_18}
\end{figure}
地球的水圈是地球全部水体及其分布的总和。地球水圈的大部分由全球性海洋构成。地球的水圈还包括大气与陆地上的水，如云、内陆海、湖泊、河流以及地下水。海洋的总质量约为 1.35×10^{18} 公吨，约为地球总质量的 1/4400。海洋覆盖面积达 3.618 亿 $km^{2}$（1.397 亿平方英里），平均深度为 3,682 m（12,080 ft），其体积估计约为 13.32 亿 $km^{3}$（3.20 亿立方英里）$^{[191]}$。

如果地球地壳表面处于同一高度，犹如一颗光滑球体，则由此形成的全球性海洋深度将为 2.7 至 2.8 km（1.68 至 1.74 mi）$^{[192]}$。约 97.5\% 的水是咸水；剩余 2.5\% 为淡水$^{[193][194]}$。大部分淡水（约 68.7\%）以冰的形式存在于冰盖和冰川中$^{[195]}$。其余 30\% 是地下水，1\% 是地表水（仅覆盖地球陆地面积的 2.8\%）$^{[196]}$，以及其他小型淡水储量，如永冻层、大气中的水汽、生物结合水等$^{[197][198]}$。

在地球最寒冷的区域，积雪在夏季仍可保存并转变为冰。这些积累的雪与冰最终形成冰川，即在自身重力作用下流动的巨大冰体。山地冰川形成于山地地区，而广阔的冰盖在极地地区的陆地上形成。冰川的流动会侵蚀地表，显著改变其形貌，形成 U 形谷及其他地貌$^{[199]}$。北极的海冰覆盖面积约与美国本土相当，然而由于气候变化，这些海冰正在迅速退缩$^{[200]}$。

地球海洋的平均盐度约为每千克海水含 35 克盐（3.5\% 的盐分）$^{[201]}$。这些盐分大多由火山活动释放，或从冷却的火成岩中析出$^{[202]}$。海洋也是大气中溶解气体的储库，这些气体对于许多水生生命形式的生存至关重要$^{[203]}$。海水对全球气候具有重要影响，海洋充当着巨大的热量储库$^{[204]}$。海洋温度分布的变化可能导致显著的气候波动，如厄尔尼诺–南方涛动现象$^{[205]}$。

地球表面丰富的水体，尤其是液态水，是使其区别于太阳系其他行星的独特特征。太阳系中具有相当大气层的行星虽部分含有大气水汽，但缺乏能够支持稳定地表液态水的表面条件$^{[206]}$。尽管一些卫星显示出存在大量地外液态水储库的迹象，其体积甚至可能超过地球海洋，但这些水体均位于数公里厚的冻结表层之下$^{[207]}$。
\subsection{大气圈}
\begin{figure}[ht]
\centering
\includegraphics[width=8cm]{./figures/07619584a1f86b23.png}
\caption{从地球的视角可以看到其大气层的不同结构：对流层中云层投下阴影，地平线上可见一条平流层的蓝色天空带，而在约 100 km 高度、太空边缘处，可见下层热层呈现的一道绿色空气辉光。} \label{fig_Earth_19}
\end{figure}
地球海平面的大气压平均为 101.325 kPa（14.696 psi）$^{[208]}$，其尺度高度约为 8.5 km（5.3 mi）$^{[3]}$。干燥大气由 78.084\% 的氮、20.946\% 的氧、0.934\% 的氩以及痕量二氧化碳和其他气体分子组成$^{[208]}$。水汽含量在 0.01\% 至 4\% 之间变化$^{[208]}$，平均约为 1\%$^{[3]}$。云层覆盖约地球表面的三分之二，其覆盖范围在海洋上多于陆地$^{[209]}$。对流层厚度随纬度变化，在两极约为 8 km（5 mi），在赤道约为 17 km（11 mi），并随天气与季节因素有所变化$^{[210]}$。

地球的生物圈显著改变了大气成分。产氧光合作用于约 27 亿年前出现，形成了当今主要由氮–氧组成的大气$^{[60]}$。这一变化使好氧生物得以繁盛，并通过大气中$O_2$随后向$O_3$的转化，间接促成了臭氧层的形成。臭氧层阻挡来自太阳的紫外辐射，使陆地生命得以存在$^{[211]}$。其他对于生命重要的大气功能包括运输水汽、提供有用气体、使小型流星在撞击地表前燃烧殆尽，以及调节温度$^{[212]}$。这一温度调节现象即温室效应：大气中的痕量分子吸收地表辐射的热能，从而提高平均温度。水汽、二氧化碳、甲烷、一氧化二氮与臭氧是主要的温室气体。若无这种热量滞留作用，地球的平均表面温度将为 −18 °C（−0.4 °F），与目前的 +15 °C（59 °F）形成鲜明对比$^{[213]}$，地球生命也很可能无法以现有形式存在$^{[214]}$。
\subsubsection{天气与气候}
\begin{figure}[ht]
\centering
\includegraphics[width=8cm]{./figures/405541a78b89eb94.png}
\caption{从太空视角看到的，分布在东太平洋及美洲上空的赤道辐合带（ITCZ）云带。} \label{fig_Earth_20}
\end{figure}
\begin{figure}[ht]
\centering
\includegraphics[width=8cm]{./figures/b083d349777c9c6a.png}
\caption{全球柯本气候分类图} \label{fig_Earth_21}
\end{figure}
地球大气没有明确的边界，而是逐渐变得稀薄并淡化进入外层空间$^{[215]}$。大气质量的四分之三集中在地表上方前 11 km（6.8 mi）内；这一最低层称为对流层$^{[216]}$。来自太阳的能量加热这一层及其下方的地表，使空气膨胀。密度较低的空气随即上升，并由更为寒冷、密度更高的空气取代。其结果是通过热能再分配形成的大气环流，从而推动天气与气候$^{[217]}$。

主要的大气环流带包括赤道区域$30^\circ$纬度以下的信风带，以及位于$30^\circ$与 $60^\circ$之间中纬度地区的西风带$^{[218]}$。海洋热含量与洋流也是决定气候的重要因素，尤其是将热能从赤道海洋输送至极地地区的温盐环流$^{[219]}$。

地球接收的太阳辐照度为 $1361 W/m^2$。$^{[3][220]}$到达地球表面的太阳能量随着纬度升高而减少。在较高纬度，阳光以更低角度照射地表，并需穿过更厚的大气柱。因此，海平面上的年平均气温自赤道起每升高一度纬度下降约 0.4 °C（0.7 °F）$^{[221]}$。地球表面可根据纬度划分为若干气候带，从赤道到极地分别为热带（或赤道带）、亚热带、温带和极地气候带$^{[222]}$。

影响某地气候的其他因素包括其与海洋的距离、海洋与大气环流以及地形$^{[223]}$。靠近海洋的地区通常夏季更凉爽、冬季更温暖，因为海洋能够储存大量热量。风将海洋的冷暖输送至陆地$^{[224]}$。大气环流也起着重要作用：旧金山与华盛顿特区均位于相近纬度的海岸，但旧金山的气候显著更加温和，因为其盛行风方向由海向陆$^{[225]}$。最后，气温随高度升高而降低，使得山地地区比低地更寒冷$^{[226]}$。

由地表蒸发产生的水汽通过大气中的环流模式被输送。当大气条件允许温暖、湿润的空气上升时，这些水汽会凝结并以降水形式落回地表$^{[217]}$。大部分水随后由河流系统输送至低海拔地区，通常回到海洋或汇入湖泊。这个水循环是支持陆地生命的重要机制，也是地质时期侵蚀地表形态的主要因素。降水分布差异极大，从每年数米的水量到不足一毫米不等。大气环流、地形特征与温度差异共同决定各区域的平均降水量$^{[227]}$。

常用的柯本气候分类系统包含五大类（湿热带、干旱带、湿润中纬度带、大陆性气候带和寒冷极地带），这些类别又可进一步细分为更具体的亚类$^{[218]}$。柯本系统依据观测到的温度与降水对区域进行分类$^{[228]}$。在诸如死亡谷等炎热沙漠中，地表气温可升至约$55^\circ C$（$131^\circ F$），而在南极洲可低至 $-89^\circ C$（$-128^\circ F$）$^{[229][230]}$。
\subsubsection{高层大气}
\begin{figure}[ht]
\centering
\includegraphics[width=8cm]{./figures/40008c3c1dad7bb5.png}
\caption{从下方向上观看地球夜侧的高层大气时，可见余辉呈带状照亮对流层，使其呈橙色，并映出云层的剪影；平流层则呈现白色与蓝色。再往上是中间层（粉色区域），其一直延伸至最低空气辉光形成的橙色与微弱绿色光带处，约位于一百千米高度，即太空边缘以及热层下缘（不可见）。更高处则可见延伸数百千米的绿色与红色的极光带。} \label{fig_Earth_22}
\end{figure}
对流层以上的大气，被统称为高层大气$^{[231]}$，通常分为平流层、中间层和热层$^{[212]}$。每一层都有不同的直减率，即温度随高度变化的速率。再往外，大气逐渐变薄，进入外逸层，最终过渡至磁层，在那里地磁场与太阳风发生相互作用$^{[232]}$。平流层内部包含臭氧层，这一成分部分屏蔽来自太阳的紫外线，因此对地球生命至关重要。卡门线位于距地表 100 km（62 mi）处，被用作大气与外层空间之间的工作定义界限$^{[233]}$。

热能使大气外缘的一些分子速度增加到足以逃离地球引力的程度。这导致了大气缓慢而持续地流失至太空。由于未固定氢的分子质量较低，它更容易达到逃逸速度，因此其泄漏速率高于其他气体$^{[234]}$。氢逸入太空促成了地球大气与地表从最初的还原性状态逐渐向当前的氧化性状态转变。光合作用提供了自由氧源，但如氢等还原剂的损失被认为是大气中氧得以广泛积累的必要前提$^{[235]}$。因此，氢能够从大气逃逸的能力可能影响了地球上生命发展的性质$^{[236]}$。在当今富氧大气中，大多数氢在逸出前即被转化为水。相反，大部分氢的损失来自上层大气中甲烷的分解$^{[237]}$。
\subsection{地球上的生命}
\begin{figure}[ht]
\centering
\includegraphics[width=8cm]{./figures/4b9c7eb24a2a51f6.png}
\caption{陆地上生产性植被密度变化的动画（棕色为低密度，深绿色为高密度），以及海洋表层浮游植物密度变化的动画（紫色为低密度，黄色为高密度）。} \label{fig_Earth_23}
\end{figure}
地球是已知唯一曾经适宜生命生存的地方。地球上的生命在地球形成后的数亿年间于早期水体中发展起来，约在 40 亿年前。地球提供了液态水——一种能够使复杂有机分子聚合与相互作用的环境——并同时提供足够的能量以维持代谢$^{[238]}$。植物和其他生物从水、土壤和大气中吸收养分，这些养分在不同物种之间不断循环$^{[239]}$。

地球上的生命塑造并居住于地球上的许多特定生态系统中，并最终扩展为全球性的、覆盖所有生态系统的生物圈$^{[240]}$。随着时间推移，地球生命也经历了巨大的多样化，使生物圈形成不同的生物群系，每个群系由相对相似的动植物所组成$^{[241]}$。不同的生物群系出现在不同的海拔高度、不同的水深、不同的行星温度纬度，以及陆地上不同的湿度条件下。地球的物种多样性和生物量在浅海区域和森林中达到峰值，尤其是在赤道附近温暖且湿润的环境中。相比之下，冰冻的极地区域与高海拔地区，或极度干旱的区域，植物与动物生命相对贫乏$^{[242]}$。

因此，生命对地球产生了影响，在漫长的时间尺度上显著改变了地球的大气与地表，造成诸如“大氧化事件”等变化$^{[243]}$。此外，人类也影响了地球、其生命以及生命的发展。
\subsubsection{地球生命所面临的挑战}
极端天气（如热带气旋）发生在地球大部分区域，并对这些地区的生命产生巨大影响。从 1980 年到 2000 年间，这类事件平均每年造成约 11,800 人死亡$^{[244]}$。许多地区还遭受地震、山体滑坡、海啸、火山喷发、龙卷风、暴风雪、洪水、干旱、野火及其他灾害$^{[245]}$。人类影响也在许多方面显现，例如空气和水污染、酸雨、植被损失（过度放牧、森林砍伐、荒漠化）、野生动物减少、物种灭绝$^{[246]}$，土壤退化、土壤枯竭与侵蚀$^{[247]}$。人类活动向大气释放温室气体，导致全球变暖$^{[248]}$。这推动了一系列变化，包括冰川和冰盖的融化、全球平均海平面上升、干旱与野火风险增加、物种向更寒冷地区迁移等$^{[249]}$。
\subsection{人文地理}
\begin{figure}[ht]
\centering
\includegraphics[width=8cm]{./figures/f6d79477027f277f.png}
\caption{一幅将地球夜间人造光源辐射叠加于地图之上的合成图像。} \label{fig_Earth_24}
\end{figure}
人类约在 30 万年前由更早的灵长类在非洲东部起源，此后不断在地球上迁徙；随着公元前第十千纪农业的出现，人类日益在地球陆地上定居$^{[250]}$。20 世纪，南极洲成为人类最后被探索与定居的大陆，尽管人类在那里的存在仍然有限。自 19 世纪以来，全球人口呈指数增长，于 2020 年代达到 80 亿$^{[251]}$，并预计将在 21 世纪下半叶达到约 100 亿的峰值$^{[252]}$。大部分增长预计将在撒哈拉以南非洲发生$^{[252]}$。

全球人类人口的分布与密度差异巨大，大多数人口居住在亚洲南部至东部地区，并有 90\% 居住在地球北半球$^{[253]}$，部分原因是世界陆地面积在北半球占优势，全球 68\% 的陆地位于北半球$^{[254]}$。此外，自 19 世纪以来，人类愈发向城市地区集中，至 21 世纪，大多数人口生活在城市中$^{[255]}$。

在地球表面之外，人类仅曾在少数特殊用途的深地及水下设施，以及少数空间站中居住。人类人口几乎完全停留在地球表面，完全依赖地球与其所维持的环境。自 20 世纪后半叶以来，已有数百人类短暂离开地球，其中仅极少数抵达另一天体——月球$^{[256][257]}$。

地球受到广泛的人类定居，人类发展出多样的社会与文化。自 19 世纪以来，地球大部分陆地被主权国家（国家）领土性宣称，并以政治边界划分，如今共存在 205 个此类国家$^{[258]}$，仅南极洲部分区域与少数小地区尚未被宣称$^{[259]}$。总体而言，这些国家大多组成联合国——全球主要的政府间国际组织$^{[260]}$——其治理范围延伸至海洋与南极洲，因此覆盖整个位地球。
\subsubsection{自然资源与土地利用}
\begin{figure}[ht]
\centering
\includegraphics[width=8cm]{./figures/300adf8b1f5261ae.png}
\caption{2019 年地球土地用于人类农业的利用状况} \label{fig_Earth_25}
\end{figure}
地球拥有被人类利用的资源$^{[261]}$。其中被称为不可再生资源的，如化石燃料，只能在地质时间尺度上得到补充$^{[262]}$。地壳中蕴藏大量化石燃料储量，包括煤、石油与天然气$^{[263]}$。这些储量被人类用于能源生产以及作为化学生产的原料$^{[264]}$。矿石矿体也通过成矿过程在地壳中形成，这一过程源于岩浆活动、侵蚀与板块构造作用$^{[265]}$。这些金属和其他元素通过采矿方式被提取，而采矿过程通常会造成环境与健康方面的损害$^{[266]}$。

地球生物圈为人类生产许多有用的生物产品，包括食物、木材、药物、氧气以及有机废物的循环利用。陆地生态系统依赖表土与淡水，而海洋生态系统依赖从陆地冲刷进入海洋的溶解养分$^{[267]}$。2019 年，地球陆地表面中有 3,900 万 $km(^2)$（1,500 万平方英里）由森林与林地构成，1,200 万 $km(^2)$（460 万平方英里）为灌木与草地，4,000 万 $km(^2)$（1,500 万平方英里）用于动物饲料生产与放牧，1,100 万 $km(^2)$（420 万平方英里）被开垦为农田$^{[268]}$。在 12–14\% 的无冰陆地农田中，2015 年有 2 个百分点为灌溉农田$^{[269]}$。人类使用天然和人造建筑材料建造居所与基础设施$^{[270]}$。
\subsubsection{人类与环境}
\begin{figure}[ht]
\centering
\includegraphics[width=8cm]{./figures/b4a619d762721997.png}
\caption{平均地表空气温度的变化及其驱动因素。人类活动导致了气温升高，而自然因素则增加了其中的变率$^{[271]}$。} \label{fig_Earth_26}
\end{figure}
人类活动已经影响了地球的环境。通过燃烧化石燃料等活动，人类增加了大气中的温室气体含量，改变了地球的能量收支和平衡，以及地球的气候$^{[248][272]}$。估计 2020 年的全球气温比工业化前基线高出 $1.2^\circ C(2.2 ^\circ F)$。$^{[273]}$这种温度升高被称为全球变暖，它促成了冰川的融化、海平面的上升、干旱与野火风险增加，以及物种向更寒冷地区迁移$^{[249]}$。

“行星边界”概念被提出以量化人类对地球的影响。在已识别的九项边界中，有五项已被跨越：生物圈完整性、气候变化、化学污染、野生栖息地破坏以及氮循环被认为已经突破安全阈值$^{[274][275]}$。截至 2018 年，没有任何国家能够在不跨越行星边界的情况下满足其人口的基本需求。然而，人们认为在可持续的资源使用水平内，全球范围仍有可能满足所有人的基本物质需求$^{[276]}$。
\subsection{文化与历史视角}
\begin{figure}[ht]
\centering
\includegraphics[width=8cm]{./figures/d3ebaa4c2d24e655.png}
\caption{2010 年 9 月 11 日，NASA 宇航员 Tracy Caldwell Dyson 在国际空间站的穹顶舱观察地球。} \label{fig_Earth_27}
\end{figure}
人类文化发展出了许多关于这颗行星的观点$^{[277]}$。地球的标准天文学符号是一个四分圆 🜨$^{[278]}$，代表世界的四个角，以及带十字的球体符号 ♁。地球有时被人格化为神祇。在许多文化中，它是一位大地母神，同时也是主要的生育神$^{[279]}$。许多宗教的创世神话都涉及由超自然神祇创造地球$^{[279]}$。20 世纪中期提出的盖娅假说（Gaia hypothesis）将地球的环境与生命视为一个单一的自我调节系统，能够实现宜居条件的广泛稳定$^{[280][281][282]}$。

从太空拍摄的地球图像，特别是在阿波罗计划期间的影像，被认为改变了人们对其所居住行星的看法，这被称为“整体观效应”（overview effect），突出强调了地球的美丽、独特性以及显而易见的脆弱性$^{[283][284]}$。尤其是这一点使人们意识到人类活动对地球环境的影响规模。在科学（尤其是地球观测技术$^{[285]}$）的推动下，人类开始在全球范围内采取行动应对环境问题$^{[286]}$，认识到人类的影响以及地球各环境之间的相互关联$^{[287]}$。

科学研究使人类对地球的认知发生了数次文化层面的重大转变。最初关于平坦地球的观念在古希腊逐渐被球形地球的概念所取代，这一观点被归功于哲学家毕达哥拉斯与巴门尼德斯$^{[288][289]}$。直到 16 世纪，人们普遍仍相信地球是宇宙的中心$^{[290]}$，直到科学家首次得出结论：地球是一个运动的天体，是太阳系行星之一$^{[291]}$。

直到 19 世纪，地质学家才意识到地球的年龄至少已有数百万年$^{[292]}$。开尔文勋爵在 1864 年使用热力学方法估算地球年龄在 2000 万至 4 亿年之间，这引发了关于该主题的激烈争论；直到 19 世纪末与 20 世纪初放射性与放射性测年被发现后，才确立了可靠的地球年龄测定机制，证明地球已有数十亿年历史$^{[293][294]}$。
\subsection{另见}
\begin{itemize}
\item 天球 —— 天文学中的概念性工具
\item 地球相 —— 从月球上观测到的地球相位变化
\item 地球科学 —— 与地球相关的自然科学领域
\item 地球上的极端现象
\item 太阳系极端现象列表
\item 地球纲要
\item 太阳系行星物理性质表
\item 地球年龄估算时间线
\item 遥远未来时间线 —— 关于极远未来的科学预测
\end{itemize}
\subsection{备注}
\begin{enumerate}
\item 所有天文数量都有所不同，两者世俗的和定期。给出的数量是瞬间的值J2000.0忽略所有的周期性变化。
\item aphelion =a× (1 +e); 近日点 =a× (1-e)，其中a是半长轴，并且e是偏心。地球的近日点和远日点之间的差异是500万公里。约翰·威尔金森 (2009)。探索新的太阳系。CSIRO出版。第144页。ISBN 978-0-643-09949-4。
\item 地球的周长几乎正好是40，000公里，因为仪表是在这个测量上校准的-更具体地说，是两极和赤道之间距离的十分之一。
\item 由于自然波动，周围的歧义冰架和映射约定垂直基准,陆地和海洋覆盖的精确值没有意义。根据来自矢量地图和全球土地覆盖 已存档2015年3月26日在回程机数据集显示，湖泊和溪流覆盖率的极值分别为地球表面的0.6\% 和1.0\%。的冰盖南极洲和格陵兰被算作陆地，尽管支撑它们的大部分岩石都位于海平面以下。
\item 最低限度的来源，[18]意思是，[19]和最大值[20]表面温度
\item 如果地球缩小到一个台球,地球上的一些地区，如大山脉和海洋海沟，会感觉像是微小的缺陷，而地球上的大部分地区，包括大平原和深渊平原,会感觉更顺畅。[88]
\item 包括索马里板块,它是由非洲板块形成的。请参阅:Chorowicz，Jean (2005年10月)。“东非裂谷系统”。非洲地球科学杂志。43(1-3):379-410。Bibcode:2005JAfES .. 43 .. 379C。doi:10.1016/j.jafrearsci.2005.07.019。
\item 局部变化之间5和200公里。
\item 局部变化之间5和70公里。
\item 这些数字的最终来源使用术语 “UT1秒” 而不是 “平均太阳时间秒”。青木，S.；木下，H.; 吉诺，B.；卡普兰，G.H、; 麦卡锡，D。D.;Seidelmann，P。K。(1982)。“世界时的新定义”。天文学和天体物理学。105(2):359-361.Bibcode:1982A & A...105 .. 359A。
\item 对于地球来说，希尔半径为
\[
R_H = a\left( \frac{m}{3M} \right)^{1/3},~
\]
其中 (m) 为地球的质量，(a) 为 1 个天文单位（AU），而 (M) 为太阳的质量。
因此，AU 尺度下的该比值的立方根约为
\[
\left( \frac{1}{3 \cdot 332{,}946} \right)^{1/3} \approx 0.01.~
\]
\item 远日点是到近日点距离的103.4\%。由于平方反比定律，近日点的辐射约为远日点能量的106.9\%。
\item 截至2018年1月4日，美国战略司令部共追踪了18，835个人造物体，其中大部分是碎片。请参阅:Anz-meador，菲利普; 射击，德比，编辑。(2018年2月)。"卫星箱得分"(PDF)。轨道碎片季刊新闻。22(1): 12。已存档(PDF)从2019年4月2日的原件。已检索4月18日2018。
\end{enumerate}
\subsection{参考资料}
\begin{enumerate}
\item Simon，J.L.; 等人 (1994年2月)。“月球和行星的进动公式和平均元素的数值表达式”。天文学和天体物理学。282(2):663-683。Bibcode:1994A & A。282 .. 663S。
\item 工作人员 (2021年3月13日)。"有用的常量"。国际地球自转和参考系统服务。已存档从原件到2012年10月29日。已检索6月8日2022年。
\item 威廉姆斯，大卫R.(2024年11月15日)。"地球概况"。NSSDCA。NASA戈达德太空飞行中心。已存档从原来的2013年5月8日。已检索12月30日2024。
\item 克拉本·沃尔特·艾伦；考克斯，亚瑟·N。(2000)。亚瑟·N.考克斯 (编)。艾伦的天体物理量。斯普林格。第294页。ISBN 978-0-387-98746-0。已存档从2023年2月21日开始。已检索3月13日2011年。
\item 各种 (2000)。大卫R.立德 (主编)。化学与物理手册(第81版)。CRC出版社。ISBN 978-0-8493-0481-1。
\item “选定的天文常数，2011年”。天文年鉴。存档自原来的2013年8月26日。已检索2月25日2011年。
\item 世界大地测量系统(WGS-84)。在线提供 已存档2020年3月11日在回程机从国家地理空间情报局。
\item 安妮·卡泽纳夫(1995)。大地水准面、地形和地貌分布(PDF)。在阿伦斯，托马斯J (ed.)。全球地球物理学: 物理常数手册。AGU参考架。第一卷。华盛顿特区: 美国地球物理联合会。Bibcode:1995年盖夫。conf ..... A。doi:10.1029/RF001。ISBN  978-0-87590-851-9。存档自原来的 (PDF)2006年10月16日。已检索8月3日2008。
\item 国际地球自转和参考系统服务 (IERS) 工作组 (2004年)。“一般定义和数字标准”(PDF)。在麦卡锡，丹尼斯D.佩蒂特，热拉尔 (主编)。IERS公约 (2003年)(PDF)。美因河畔法兰克福: 联邦法院。第12页。ISBN  978-3-89888-884-4。已存档 (PDF)从原来的2016年8月12日。已检索4月29日2016。
\item Humerfelt，Sigurd (2010年10月26日)。“WGS 84如何定义地球”。家庭在线。存档自原来的2011年4月24日。已检索4月29日2011年。
\item 迈克尔·皮德维尼 (2006年2月2日)。“8(o)。海洋简介 | 表8o-1: 我们星球被海洋和大陆覆盖的表面积。。自然地理学基础，第2版。不列颠哥伦比亚大学，奥肯那根。已存档从2006年12月9日的原件。已检索11月26日2007。
\item “行星物理参数”。喷气推进实验室。2008。已检索8月11日2022年。
\item 国际单位制 (SI)(PDF)(2008年版)。美国商务部,NIST特别出版物330。第52页。存档自原来的(PDF)2009年2月5日
\item 威廉姆斯，詹姆斯G.(1994)。“对地球倾角、进动和章动的贡献”。天文杂志。108: 711。Bibcode:1994AJ....108..711W。doi:10.1086/117108。ISSN0004-6256。S2CID122370108。  
\item 克拉本·沃尔特·艾伦；考克斯，亚瑟·N。(2000)。亚瑟·N.考克斯 (编)。艾伦的天体物理量。斯普林格。第296页。ISBN 978-0-387-98746-0。已存档从2023年2月21日开始。已检索8月17日2010。
\item 克拉本·沃尔特·艾伦；考克斯，亚瑟·N。(2000)。亚瑟·N.考克斯 (编)。艾伦的天体物理量(第四版)。纽约: AIP出版社。第244页。ISBN 978-0-387-98746-0。已存档从2023年2月21日开始。已检索8月17日2010。
\item “大气和行星温度”。美国化学学会。2013年7月18日存档自原来的2023年1月27日。已检索1月3日2023。
\item “世界: 最低温度”。气象组织极端天气和气候档案。亚利桑那州立大学。已存档从原件到2019年3月23日。已检索9月6日2020。
\item 琼斯，P。D.；哈珀姆，C。(2013)。“地球绝对表面气温的估计”。地球物理研究杂志: 大气。118(8):3213-3217。Bibcode:2013JGRD..118.3213J。doi:10.1002/jgrd.50359。ISSN2169-8996。 
\item “世界: 最高温度”。气象组织极端天气和气候档案。亚利桑那州立大学。已存档从2018年5月1日开始。已检索9月6日2020。
\item 联合国原子辐射影响问题科学委员会 (2008年)。电离辐射的来源和影响。纽约: 联合国 (2010年出版)。表1。ISBN 978-92-1-142274-0。已存档从原始的2019年7月16日。已检索11月9日2012。
\item “地球，n.1”。牛津英语词典(第3版)。牛津,英格兰:牛津大学出版社。2010。doi:10.1093/acref/9780199571123.001.0001。ISBN 978-0-19-957112-3。
\item 鲁道夫·西梅克(2007)。北方神话词典。由安吉拉霍尔翻译。D.S。布鲁尔。第179页。ISBN 978-0-85991-513-7。
\item “地球”。新牛津英语词典(第一版)。牛津:牛津大学出版社。1998。ISBN 978-0-19-861263-6。
\item "Terra"。牛津英语词典(在线版)。牛津大学出版社。 (订阅或参与机构成员必需。)
\item "Tellus"。牛津英语词典(在线版)。牛津大学出版社。 (订阅或参与机构成员必需。)
\item "盖亚"。牛津英语词典(在线版)。牛津大学出版社。 (订阅或参与机构成员必需。)
\item "Terran"。牛津英语词典(在线版)。牛津大学出版社。 (订阅或参与机构成员必需。)
\item "地面"。牛津英语词典(在线版)。牛津大学出版社。 (订阅或参与机构成员必需。)
\item "terrene"。牛津英语词典(在线版)。牛津大学出版社。 (订阅或参与机构成员必需。)
\item "碲化物"。牛津英语词典(在线版)。牛津大学出版社。 (订阅或参与机构成员必需。)
\item "碲"。Lexico英国英语词典。牛津大学出版社。存档自原来的2021年3月31日。
\item 奥黛丽·布维尔；Wadhwa, Meenakshi(2010年9月)。“太阳系的年龄由陨石包裹体中最古老的pb-pb年龄重新定义”。自然地球科学。3(9):637-641。Bibcode:2010年自然。3。637B。doi:10.1038/ngeo941。
\item 请参阅:
\begin{itemize}
\item 达尔林普尔，G。布伦特(1991)。地球的年龄。加州: 斯坦福大学出版社。ISBN 978-0-8047-1569-0。
\item 纽曼，威廉L.(2007年7月9日)。《地球时代》。出版物服务，美国地质勘探局。已存档从2005年12月23日的原件开始。已检索9月20日2007。
\item 达尔林普尔，G。布伦特(2001)。“二十世纪的地球时代: 一个 (大部分) 解决的问题”。地质学会，伦敦，特别出版物。190(1):205-221.Bibcode:2001GSLSP.190..205D。doi:10.1144/GSL.SP.2001.190.01.14。S2CID 130092094。已存档从2007年11月11日的原件。已检索9月20日2007。
\end{itemize}
\item 右特，K.；Schonbachler, M.(2018年5月7日)。“吸积期地幔的Ag同位素演化: Pd和Ag金属硅酸盐分配的新约束”。差异化: 构建行星的内部架构。2084: 4034。Bibcode:2018LPICo2084.4034R。已存档从原件到2020年11月6日。已检索10月25日2020。
\item 塔特塞，罗曼; 阿纳德，马赫什; 加塔切卡，杰罗姆；乔伊，凯瑟琳·H。；莫蒂默，詹姆斯一世；佩内特-费舍尔，约翰·F；萨拉·罗素；斯内普，约书亚·F；韦斯，本杰明·P。(2019年)。“限制月球和内部太阳系的进化历史: 新返回的月球样本的情况”。空间科学评论。215(8): 54。Bibcode:2019SSRv .. 215...54T。doi:10.1007/s11214-019-0622-x。ISSN1572-9672。 
\item 赖利，迈克尔 (2009年10月22日)。“有争议的月球起源理论改写历史”。发现新闻。存档自原来的2010年1月9日。已检索1月30日2010。
\item Canup, R.;阿斯霍格，E.I.(2001)。“月球的起源在地球形成的尽头附近发生了巨大的撞击”。自然。412(6848):708-712。Bibcode:2001 Natur.412..708C。doi:10.1038/35089010。PMID11507633。S2CID4413525。  
\item 梅尔，M.M。M、；Reufer, A.;维勒，R。(2014年8月4日)。“关于Theia的起源和组成: 来自巨大冲击新模型的约束”。伊卡洛斯。242: 5。arXiv:1410.3819。Bibcode:2014年Icar .. 242 .. 316米。doi:10.1016/j.icarus.2014.08.003。ISSN0019-1035。S2CID119226112。  
\item 克莱斯，菲利普；亚历山德罗·莫比德利(2011)。“晚轰炸”。穆里尔在加尔高德; 阿米尔斯，里卡多教授; 金塔尼利亚，何塞·塞尼查罗; 克莱夫二世，亨德森·詹姆斯 (吉姆); 欧文，威廉·M；平提，丹尼尔·L教授；米歇尔·维索 (编辑)。天体生物学百科全书。施普林格柏林海德堡.pp. 909-912。doi:10.1007/978-3-642-11274-4_869。ISBN 978-3-642-11271-3。
\item 地球早期的大气层和海洋。月球和行星研究所。大学空间研究协会。已存档从原始的2019年7月8日。已检索6月27日2019。
\item 莫比德利，A.等 (2000年)。“向地球输送水的来源区域和时间尺度”。气象学与行星科学。35(6):1309-1320。Bibcode:2000M & PS。35.1309米。doi:10.1111/j.1945-5100.2000.tb01518.x。
\item 皮亚尼，劳雷特; 等人 (2020年)。“地球上的水可能是从类似于顽辉石球粒陨石的物质中继承而来的”。科学。369(6507):1110-1113。Bibcode:2020Sci。369.1110页。doi:10.1126/科学.aba1948。ISSN0036-8075。PMID32855337。S2CID221342529。   
\item 桂南，E.F.;里巴斯，我。(2002)。本杰明·蒙特西诺斯、阿尔瓦罗·吉梅内斯和爱德华·F。桂南 (编)。我们不断变化的太阳: 太阳核演化和磁活动对地球大气和气候的作用。ASP会议论文集: 不断发展的太阳及其对行星环境的影响。旧金山: 太平洋天文学会。Bibcode:2002ASPC..269...85G。ISBN 978-1-58381-109-2。
\item 工作人员 (2010年3月4日)。地球磁场的最古老的测量揭示了太阳和地球之间的战斗，我们的大气层。Phys.org。已存档从2011年4月27日的原始。已检索3月27日2010。
\item 培训师，Melissa G.; 等人 (2006年11月28日)。“土卫六和早期地球上的有机雾霾”。美国国家科学院院刊。103(48):18035-18042。doi:10.1073/pnas.0608561103。ISSN0027-8424。PMC1838702。PMID17101962。   
\item 麦克多诺，W.F.；太阳，S.-s.(1995)。“地球的组成”。化学地质学。120(3-4):223-253。Bibcode:1995 chgeo.120..223米。doi:10.1016/0009-2541(94)00140-4。已存档从2023年5月6日开始。已检索5月6日2023。
\item 哈里森，T。M。;Blichert-Toft, J.；穆勒，W.；Albarede, F.；霍尔顿，P.；Mojzsis, S.(2005年12月)。“异质哈德铪: 4.4至4.5 ga的大陆地壳的证据”。科学。310(5756):1947年-1950年。Bibcode:2005Sci。310-1947小时。doi:10.1126/science.1117926。PMID16293721。S2CID11208727。  
\item 罗杰斯，约翰·詹姆斯·威廉; 桑托什，M。(2004)。大洲和超大陆。牛津大学出版社美国。第48页。ISBN 978-0-19-516589-0。
\item 赫尔利，P。M、；兰德，J.R。(1969年6月)。“漂移前大陆核”。科学。164(3885):1229-1242。Bibcode:1969Sci。164.1229小时。doi:10.1126/science.164.3885.1229。PMID17772560。 
\item 阿姆斯特朗，R。L.(1991)。“地壳增长的持续神话”(PDF)。澳大利亚地球科学杂志。38(5):613-630。Bibcode:1991年奥耶斯。。38 .. 613A。CiteSeerX10.1.1.527.9577。doi:10.1080/08120099108727995。已存档(PDF)从原始的2017年8月8日。已检索10月24日2017。  
\item De Smet, J.;范登伯格，a.p.；Vlaar, N.J.(2000)。“对流地幔减压融化导致的大陆的早期形成和长期稳定性”(PDF)。构造物理学。322(1-2):19-33。Bibcode:2000 Tectp.322. .. 19D。doi:10.1016/S0040-1951(00)00055-x。高密度脂蛋白:1874/1653。已存档从原件到2021年3月31日。已检索8月25日2019。
\item Dhuime, B.;霍克斯沃斯，C.J.；Delavault，H.; 卡伍德，P.A.(2018)。“大陆壳的生成和破坏速率: 对大陆增长的影响”。哲学交易A。376(2132)。Bibcode:2018RSPTA.37670403D。doi:10.1098/rsta.2017.0403。PMC6189557。PMID30275156。  
\item 布拉德利，华盛顿特区.(2011)。“地质记录和超大陆周期的长期趋势”。地球科学评论。108(1-2):16-33。Bibcode:2011ESRv .. 108...16B。CiteSeerX10.1.1.715.6618。doi:10.1016/j.earscirev.2011.05.003。S2CID140601854。  
\item Kinzler, Ro.“冰河时代何时以及如何结束？“另一个可以开始吗？”。学。美国自然历史博物馆。已存档从原件到2019年6月27日。已检索6月27日2019。
\item Chalk，Thomas B.; 等人 (2007年12月12日)。“中更新世过渡时期冰河时代加剧的原因”。科学。114(50):13114-13119。doi:10.1073/pnas.1702143114。PMC5740680。PMID29180424。  
\item 工作人员。“古气候学-古代气候的研究”。Page古生物学科学中心。存档自原来的2007年3月4日。已检索3月2日2007。
\item 特纳，克里斯·S·M; 等人 (2010年)。“新西兰贝壳杉 (Agathis australis) 在最后一个冰期间隔 (60，000-11，700年前) 测试突然气候变化的同步性的潜力”。第四纪科学评论。29(27-28)。爱思唯尔:3677-3682。Bibcode:2010QSRv。29.3677t。doi:10.1016/j.quascirev.2010.08.017。已存档从原件到2021年3月31日。已检索11月3日2020。
\item 杜立特尔，W.福特;蠕虫，鲍里斯(2000年2月)。“连根拔起生命之树”(PDF)。科学美国人。282(6):90-95。Bibcode:2000 SciAm.282b .. 90D。doi:10.1038/scientificamerican0200-90。PMID10710791。存档自原来的(PDF)2011年7月15日  
\item 卡尔·齐默(2013年10月3日)。“地球的氧气: 一个容易被视为理所当然的谜团”。纽约时报。存档自原来的2013年10月3日。已检索10月3日2013。
\item 伯克纳，L。V。；马歇尔，L。C.(1965)。“关于地球大气中氧气浓度的起源和上升”。大气科学杂志。22(3):225-261.Bibcode:1965年JAtS...22 .. 225B。doi:10.1175/1520-0469(1965)022<0225:OTOARO>2.0.CO;2。
\item 凯瑟琳·伯顿 (2002年11月29日)。“天体生物学家发现了陆地上早期生命的证据”。NASA.存档自原来的2011年10月11日。已检索3月5日2007。
\item 诺夫克，诺拉；克里斯蒂安，丹尼尔; 韦西，大卫；哈森，罗伯特·M。(2013年11月8日)。“微生物诱导的沉积结构记录了西澳大利亚州皮尔巴拉约34.8亿年历史的Dresser地层中的一个古老生态系统”。天体生物学。13(12):1103-1124。Bibcode:2013阿斯比奥。。13.1103N。doi:10.1089/ast.2013.1030。PMC3870916。PMID24205812。  
\item Ohtomo，Yoko; Kakegawa，Takeshi; Ishida，Akizumi; 等人 (2014年1月)。“早期太古宙Isua沉积沉积岩中生物成因石墨的证据”。自然地球科学。7(1):25-28.Bibcode:2014NatGe...7...25O。doi:10.1038/ngeo2025。ISSN1752-0894。S2CID54767854。  
\item Borenstein，Seth (2015年10月19日)。“在被认为是荒凉的早期地球上的生命暗示”。激发。纽约扬克斯:Mindspark互动网络。美联社。存档自原来的2016年8月18日。已检索10月20日2015。
\item 伊丽莎白·贝尔；博赫尼克，帕特里克；哈里森，T。标记;毛，温迪L。(2015年10月19日)。“保存在41亿年历史的锆石中的潜在生物碳”。Proc.Natl.Acad。Sci。美国。112(47):14518-4521。Bibcode:2015PNAS..11214518B。doi:10.1073/pnas.1517557112。ISSN1091-6490。PMC4664351。PMID26483481。   早期版本，印刷前在线出版。
\item Tyrell，Kelly April (2017年12月18日)。“迄今为止发现的最古老的化石表明地球上的生命始于35亿年前”。威斯康星大学麦迪逊分校。已存档从原件到2021年3月31日。已检索12月18日2017。
\item Schopf, J.威廉; 北岛，Kouki; Spicuzza，Michael J.；库德里亚夫采夫，阿纳托利B.；谷，约翰·W。(2017)。“SIMS对已知最古老的微化石组合的分析记录了它们与分类群相关的碳同位素组成”。PNAS。115(1):53-58。Bibcode:2018PNAS..115...53S。doi:10.1073/pnas.1718063115。PMC5776830。PMID29255053。  
\item “地球-月球动力学”。月球和行星研究所。已存档从原来的2015年9月7日。已检索9月2日2022年。
\item 布鲁克，约翰L.(2014)。气候变化与全球历史进程。剑桥大学出版社。第42页。ISBN 978-0-521-87164-8。
\item Cabej, Nelson R.(2019年)。寒武纪大爆发的表观遗传机制。爱思唯尔科学。第56页。ISBN 978-0-12-814312-4。
\item 斯坦利，S.M。(2016)。“估计地球历史上主要海洋物种大灭绝的规模”。美国国家科学院院刊。113(42):E6325-E6334。Bibcode:2016PNAS..113E6325S。doi:10.1073/pnas.1613094113。PMC5081622。PMID27698119。S2CID23599425。    
\item 古尔德，斯蒂芬J.(1994年10月)。“地球上生命的进化”。科学美国人。271(4):84-91。Bibcode:1994 SciAm.271d .. 84克。doi:10.1038/scientificamerican1094-84。PMID7939569。已存档从2007年2月25日的原件。已检索3月5日2007。 
\item 达弗，G.；盖伊，F.；麦凯，H. T.；Likius, A.;Boisserie，J.-R；穆萨，A.；帕拉斯，L.；Vignaud, P.;Clarisse, N.D.(2022年)。“乍得中新世晚期人类双足主义的颅后证据”。自然。609(7925):94-100。Bibcode:2022年自然609. .. 94D。doi:10.1038/s41586-022-04901-z。ISSN1476-4687。PMID36002567。已存档从2022年8月27日开始。已检索3月29日2024。  
\item 威尔金森，B。H、; 麦克罗伊，B。J.(2007)。“人类对大陆侵蚀和沉积的影响”。美国地质学会公报。119(1-2):140-156。Bibcode:2007GSAB..119..140W。doi:10.1130/B25899.1。S2CID128776283。 
\item 萨克曼，I.-J；Boothroyd, A.I.;克莱默，K。E。(1993)。“我们的太阳.III.“现在和未来”。天体物理学杂志。418:457-468.Bibcode:1993年apj。418 .. 457S。doi:10.1086/173407。
\item 罗伯特·布里特 (2000年2月25日)。“冻结，油炸或干燥: 地球有多长时间？”。Space.com。存档自原来的2009年6月5日
\item 李金辉; 帕莱万，卡维；Kirschvink, Joseph L.;容，Yuk L.(2009)。“大气压力作为具有生物圈的地球行星的自然气候调节器”(PDF)。美国国家科学院院刊。106(24):9576-9579。Bibcode:2009PNAS..106.9576L。doi:10.1073/pnas.0809436106。PMC2701016。PMID19487662。已存档(PDF)从2009年7月4日开始。已检索7月19日2009年。   
\item 沃德，彼得D.;唐纳德·布朗利(2002)。地球的生与死: 天体生物学的新科学如何绘制我们世界的最终命运。纽约: 时代书籍，亨利·霍尔特和公司。ISBN 978-0-8050-6781-1。
\item 梅洛，费尔南多·德索萨; 弗里亚萨，阿蒙西奥·塞萨尔·桑托斯 (2020年)。“地球上生命的终结不是世界末日: 收敛到生物圈寿命的估计？”。国际天体生物学杂志。19(1):25-42。Bibcode:2020年IJAsB .. 19。25D。doi:10.1017/S1473550419000120。ISSN1473-5504。 
\item 布纳马，克里斯汀; 弗兰克，S.；冯·布洛，W。(2001)。“地球海洋的命运”。水文学与地球系统科学。5(4):569-575。Bibcode:2001HESS....5..569B。doi:10.5194/hess-5-569-2001。S2CID14024675。 
\item 施罗德，K.-P；康农·史密斯，罗伯特 (2008)。“太阳和地球的遥远未来重新审视”。皇家天文学会月刊。386(1):155-163.arXiv:0801.4031。Bibcode:2008MNRAS.386..155S。doi:10.1111/j.1365-2966.2008.13022.x。S2CID10073988。 
另请参见杰森·帕尔默 (2008年2月22日)。“希望地球能在太阳的死亡中幸存下来”。NewScientist.com新闻服务。存档自原来的2012年4月15日。已检索3月24日2008。
\item 霍纳，琼蒂 (2021年7月16日)。“我一直想知道: 为什么恒星，行星和卫星是圆的，而彗星和小行星不是？”。对话。已存档从原始的2023年3月3日。已检索3月3日2023。
\item Lea，罗伯特 (2021年7月6日)。“地球有多大？”。Space.com。存档自原来的2024年1月9日。已检索1月11日2024。
\item Sandwell, D.T.;史密斯，沃尔特H. F.(2006年7月7日)。“用卫星高度计数据探索海洋盆地”。NOAA/NGDC。存档自原来的2014年7月15日。已检索4月21日2007。
\item 米尔伯特，D.G.;史密斯，D。A.“使用GEOID96大地水准面高度模型将GPS高度转换为NAVD88高程”。国家大地测量局。已存档从2011年8月20日的原件。已检索3月7日2007。
\item 希瑟·A·斯图尔特；杰米森，艾伦J.(2019年)。“五深: 世界上每个海洋中最深的地方的位置和深度”。地球科学评论。197102896。Bibcode:2019ESRv .. 19702896S。doi:10.1016/j.earscirev.2019.102896。ISSN0012-8252。 
\item “撞球比地球更光滑吗？”(PDF)。台球文摘.2013年6月1日。已存档(PDF)从原来的2014年9月4日。已检索11月26日2014。
\item 芭芭拉·特克斯伯里。“信封的计算: 喜马拉雅山的规模”。卡尔顿大学。已存档从原件到2020年10月23日。已检索10月19日2020。
\item Senne，Joseph H. (2000)。“埃德蒙·希拉里爬错了山”。专业测量师。20(5):16-21。存档自原来的2015年7月17日。已检索7月16日2015。
\item 罗伯特·克鲁维奇(2007年4月7日)。“地球上 '最高' 的地方”。NPR.org。已存档从原来的2013年1月30日。已检索7月31日2012。
\item “海洋表面地形”。来自太空的海洋表面地形。NASA。已存档从原件到2021年7月29日。已检索6月16日2022年。
\item “大地水准面是什么？”。oceanservice.noaa.gov。国家海洋服务。已存档从原件到2020年10月17日。已检索10月10日2020。
\item 贾宁，H.; 曼迪亚，S.A.(2012)。海平面上升: 原因和影响简介。麦克法兰公司，出版社。第20页。ISBN 978-0-7864-5956-8。已存档从2023年2月21日开始。已检索8月26日2022年。
\item 罗，克里斯汀 (2020年2月3日)。“是海洋还是海洋？”。福布斯。已存档从2022年8月26日开始。已检索8月26日2022年。
\item 史密斯，伊维特 (2021年6月7日)。“地球是一个水世界”。NASA。已存档从2022年8月27日开始。已检索8月27日2022年。
\item “水-世界”。国家地理学会。2022年5月20日。已存档从2022年8月19日开始。已检索8月24日2022年。
\item 鲁宁，乔纳森一世。(2017)。“海洋世界探索”。宇航学报。131。Elsevier BV:123-130。Bibcode:2017 AcAau.131..123L。doi:10.1016/j.actaastro.2016.11.017。ISSN0094-5765。 
\item “海洋世界”。海洋世界。存档自原来的2022年8月27日。已检索8月27日2022年。
\item 保罗·沃森 (2021年3月9日)。“古代地球是一个水世界。科学。371(6534)。美国科学促进协会 (AAAS):1088-1089。doi:10.1126/科学.abh4289。ISSN0036-8075。PMID33707245。S2CID241687784。   
\item “NOAA海洋探险家: 加拉帕戈斯群岛: 山脊与热点相遇的地方”。oceanexplorer.noaa.gov。存档自原来的2023年11月15日。已检索4月28日2024。
\item 罗斯·邓恩；米切尔，劳拉·J；沃德，克里 (2016)。新世界历史: 教师和研究人员的实地指南。加州大学出版社。第232页-。ISBN 978-0-520-28989-5。已存档从2023年2月21日开始。已检索8月9日2023。
\item Dempsey，Caitlin (2013年10月15日)。“关于世界各大洲的地理事实”。地理领域。已存档从2022年8月26日开始。已检索8月26日2022年。
\item R.W.McColl，ed. (2005)。“大陆”。世界地理百科全书。第一卷。档案事实，公司，第215页。ISBN 978-0-8160-7229-3。已存档从2023年2月21日开始。已检索8月25日2022年。由于非洲和亚洲在苏伊士半岛相连，欧洲，非洲和亚洲有时被合并为非洲-欧亚大陆或欧亚大陆。国际奥林匹克委员会的官方旗帜，包含美洲的单一大陆 (北美和南美连接为巴拿马地峡)。
\item 国家地球物理数据中心 (2020年8月19日)。“来自依托泊1的地球表面的Hypsographic曲线”。ngdc.noaa.gov。已存档从2017年9月15日的原件。已检索9月15日2017。
\item 卡洛维茨，迈克尔; 西蒙，罗伯特 (2019年7月15日)。“为树木和碳看到森林: 在三个维度上绘制世界森林”。NASA地球观测站。已存档从2022年12月31日开始。已检索12月31日2022年。
\item “冰原”。国家地理学会。2006年8月6日。已存档从2023年11月27日开始。已检索1月3日2023。
\item Obu, J.(2021)。“地球表面有多少被永久冻土覆盖？”地球物理研究杂志: 地球表面。126(5) e2021JF006123。美国地球物理联合会 (AGU)。Bibcode:2021JGRF..12606123O。doi:10.1029/2021jf006123。ISSN2169-9003。S2CID235532921。  
\item 凯恩，弗雷泽 (2010年6月1日)。“地球表面有多少百分比是沙漠？”。今天的宇宙。已存档从原始的2023年1月3日。已检索1月3日2023。
\item “世界银行耕地”。世界银行。已存档从原来的2015年10月2日。已检索10月19日2015。
\item “世界银行永久耕地”。世界银行。已存档从原来的2015年7月13日。已检索10月19日2015。
\item 胡克，罗杰·莱布；何塞·F·马丁·杜克；Pedraza，Javier (2012年12月)。“人类的土地改造: 回顾”(PDF)。GSA今天。22(12):4-10。Bibcode:2012GSAT...12l...4H。doi:10.1130/GSAT151A.1。已存档(PDF)从2018年1月9日开始。已检索1月9日2018。
\item 工作人员。“地球的层次”。火山世界。俄勒冈州立大学存档自原来的2013年2月11日。已检索3月11日2007。
\item 杰西，大卫。“风化和沉积岩”。加州州立理工大学波莫纳分校。存档自原来的2007年7月3日。已检索3月20日2007。
\item Kring, David A.“陆地撞击陨石坑及其环境影响”。月球和行星实验室。已存档从2011年5月13日的原件。已检索3月22日2007。
\item 马丁，罗纳德 (2011)。地球的进化系统: 地球的历史。琼斯和巴特利特学习。ISBN 978-0-7637-8001-2。OCLC 635476788。已存档从2023年2月21日开始。已检索8月9日2023。
\item 布朗，W.K.;Wohletz, K.H. (2005年)。“SFT和地球的构造板块”。洛斯阿拉莫斯国家实验室。已存档从2016年4月2日的原件。已检索3月2日2007。
\item Kious, W.J.;Tilling, R.I.(1999年5月5日)。“了解板块运动”。美国地质勘探局。已存档从2011年8月10日的原件。已检索3月2日2007。
\item 考特尼·塞利格曼 (2008)。“类地行星的结构”。在线天文学eText目录。cseligman.com。已存档从2008年3月22日的原件。已检索2月28日2008。
\item Duennebier，Fred (1999年8月12日)。"太平洋板块运动"。夏威夷大学。已存档从2011年8月31日的原件。已检索3月14日2007。
\item 穆勒，R。D.; 等人 (2007年3月7日)。《海底时代》海报。NOAA。已存档从2011年8月5日的原件。已检索3月14日2007。
\item 鲍林，塞缪尔·A。威廉姆斯，伊恩·S。(1999)。“来自加拿大西北部的普里斯科安 (4.00-4.03 Ga) 正交系”。对矿物学和岩石学的贡献。134(1):3-16.Bibcode:1999年比较。。134。3B。doi:10.1007/s004100050465。S2CID128376754。 
\item Meschede，Martin; Barckhausen，Udo (2000年11月20日)。“科科斯-纳斯卡传播中心的板块构造演化”。海洋钻探计划论文集。德克萨斯A & M大学已存档从2011年8月8日的原件。已检索4月2日2007。
\item 阿格斯，d.f.；戈登，r.g.；德梅茨，C。(2011)。“56个板块相对于无净旋转参考系的地质当前运动”。地球化学，地球物理学，地球系统。12(11): 不适用。Bibcode:2011GGG....1211001A。doi:10.1029/2011GC003751。
\item 乔丹，T。H. (1979)。“地球内部的结构地质学”。美国国家科学院院刊。76(9):4192-4200。Bibcode:1979PNAS...76.4192J。doi:10.1073/pnas.76.9.4192。PMC411539。PMID16592703。  
\item 尤金·C·罗伯逊.(2001年7月26日)。“地球的内部”。美国地质勘探局。已存档从2011年8月28日的原件。已检索3月24日2007。
\item “地壳和岩石圈”。伦敦地质学会。2012。已存档从原件到2020年10月28日。已检索10月25日2020。
\item 米卡利齐奥，卡里尔-苏; 埃弗斯，珍妮 (2015年5月20日)。“岩石圈”。国家地理。已存档从2022年5月29日开始。已检索10月13日2020。
\item 谷本俊郎 (1995)。“地球的地壳结构”(PDF)。在托马斯J。阿伦斯 (主编)。全球地球物理学: 物理常数手册。AGU参考架。第一卷。华盛顿特区: 美国地球物理联合会。Bibcode:1995年盖夫。conf ..... A。doi:10.1029/RF001。ISBN  978-0-87590-851-9。存档自原来的 (PDF)2006年10月16日。已检索2月3日2007。
\item Deuss，Arwen (2014)。“地球内核的异质性和各向异性”。Annu。Rev.地球星球。Sci。42(1):103-126。Bibcode:2014AREPS..42..103D。doi:10.1146/annurev-earth-060313-054658。已存档从原件到2020年5月7日。已检索2月8日2023。
\item 摩根，J。W.;安德斯，E.(1980)。地球、金星和水星的化学成分。美国国家科学院院刊。77(12):6973-6977。Bibcode:1980PNAS...77.6973M。doi:10.1073/pnas.77.12.6973。PMC350422。PMID16592930。  
\item 布朗，杰夫·C；艾伦·穆塞特.(1981)。人迹罕至的地球(第二版)。泰勒和弗朗西斯。p. 166。ISBN 978-0-04-550028-4。注: 在Ronov和Yaroshevsky (1969) 之后。
\item 唐纳德·L.特科特; 杰拉尔德·舒伯特 (2002年3月25日)。地球动力学。剑桥大学出版社。ISBN 978-0-521-66624-4。
\item 罗伯特·桑德斯 (2003年12月10日)。放射性钾可能是地球核心的主要热源。加州大学伯克利分校新闻。已存档从原来的2013年8月26日。已检索2月28日2007。
\item “地球中心比以前想象的要热1000度”。欧洲同步加速器 (ESRF)。2013年4月25日。存档自原来的2013年6月28日。已检索4月12日2015。
\item 阿尔菲，D.；吉兰，M.J.;Vočadlo, L.;布罗德霍特，J.；价格，G。D.(2002)。"The从头算“地球核心的模拟”(PDF)。皇家学会的哲学交易。360(1795):1227-1244。Bibcode:2002RSPTA.360.1227A。doi:10.1098/rsta.2002.0992。PMID12804276。S2CID21132433。已存档(PDF)从2009年9月30日的原件。已检索2月28日2007。   
\item 特科特，D。L.；舒伯特，G。(2002)。“4”。地球动力学(2版)。英国剑桥: 剑桥大学出版社。第137页。ISBN 978-0-521-66624-4。
\item Vlaar, N; Vankeken, P.;范登堡，A.(1994)。“太古宙地球的冷却: 在较热的地幔中压力释放融化的后果”(PDF)。地球和行星科学快报。121(1-2):1-18.Bibcode:1994E & PSL.121....1V。doi:10.1016/0012-821X(94)90028-0。存档自原来的(PDF)2012年3月19日
\item 波拉克，亨利·N.；苏珊娜·J·赫特；约翰逊，杰弗里R.(1993年8月)。“来自地球内部的热流: 全球数据集分析”。地球物理学的评论。31(3):267-280。Bibcode:1993年RvGeo .. 31 .. 267页。doi:10.1029/93RG01249。
\item 理查兹，M。A.;邓肯，R。A.;Courtillot，V。E。(1989)。“洪水玄武岩和热点轨道: 羽流的头部和尾部”。科学。246(4926):103-107.Bibcode:1989Sci...246 .. 103R。doi:10.1126/science.246.4926.103。PMID17837768。S2CID9147772。  
\item 约翰·G·斯克拉特; 巴里·帕森斯；克劳德·乔皮尔(1981)。“海洋和大陆: 热量损失机制的异同”。地球物理研究杂志。86(B12) 11535。Bibcode:1981JGR....8611535S。doi:10.1029/JB086iB12p11535。
\item 瓦茨，A.B.;戴利，S.F.(1981年5月)。“长波重力和地形异常”。地球与行星科学年度评论。9(1):415-418。Bibcode:1981AREPS...9..415W。doi:10.1146/annurev.ea.09.050181.002215。
\item 彼得·奥尔森; 哈盖·阿米特 (2006年)。“地球偶极子的变化”(PDF)。Naturwissenschaften。93(11):519-542。Bibcode:2006NW.....93..519O。doi:10.1007/s00114-006-0138-6。PMID16915369。S2CID22283432。已存档(PDF)从原件到2019年9月27日。已检索7月6日2019。   
\item 理查德·菲茨帕特里克 (2006年2月16日)。“MHD发电机理论”。NASA WMAP。已存档从原件到2020年4月27日。已检索2月27日2007。
\item 坎贝尔，华莱士·霍尔 (2003)。地磁场导论。纽约: 剑桥大学出版社。第57页。ISBN 978-0-521-82206-0。
\item Ganushkina, N.Yu; Liemohn，M。W.;杜比亚金，S.(2018)。“地球磁层中的当前系统”。地球物理学的评论。56(2):309-332。Bibcode:2018年RvGeo .. 56 .. 309克。doi:10.1002/2017RG000590。高密度脂蛋白:2027.42/145256。ISSN1944-9208。S2CID134666611。存档自原来的2021年3月31日。已检索10月24日2020。  
\item Masson，Arnaud (2007年5月11日)。“星系团揭示了地球弓形冲击的改革”。欧洲航天局。已存档从原件到2021年3月31日。已检索8月16日2016。
\item 加拉格尔，丹尼斯L.(2015年8月14日)。“地球的等离子层”。NASA/马歇尔太空飞行中心。已存档从原来的2016年8月28日。已检索8月16日2016。
\item 加拉格尔，丹尼斯L.(2015年5月27日)。“等离子层是如何形成的”。NASA/马歇尔太空飞行中心。存档自原来的2016年11月15日。已检索8月16日2016。
\item 鲍姆约翰，沃尔夫冈; 特鲁曼，鲁道夫A。(1997)。基本空间等离子体物理学。《世界科学》，第8、31页。ISBN 978-1-86094-079-8。
\item 麦克罗伊，迈克尔·B。(2012)。“电离层和磁层”。大英百科全书。大英百科全书.已存档从原来的2016年7月3日。已检索8月16日2016。
\item 范艾伦，詹姆斯·阿尔弗雷德 (2004)。磁层物理学的起源。爱荷华大学出版社。ISBN 978-0-87745-921-7。OCLC 646887856。
\item 斯特恩，大卫·P。(2005年7月8日)。“地球磁层的探索”。NASA.存档自原来的2013年2月14日。已检索3月21日2007。
\item 麦卡锡，丹尼斯D.；哈克曼，克里斯汀; 纳尔逊，罗伯特·A。(2008年11月)。《闰秒的物理基础》。天文杂志。136(5):1906年-1908年。Bibcode:2008AJ....136.1906M。doi:10.1088/0004-6256/136/5/1906。
\item "闰秒"。时间服务部，美国编号。存档自原来的2015年3月12日。已检索9月23日2008。
\item “地球方向的快速服务/预测”。IERS公告-A。28(15)。2015年4月9日存档自原来的(.DAT文件 (在浏览器中显示为明文))2015年3月14日。已检索4月12日2015。
\item Seidelmann，P。肯尼斯 (1992)。天文历书的解释性补充。米尔谷，CA: 大学科学书籍。第48页。ISBN 978-0-935702-68-2。
\item Zeilik，迈克尔; 格雷戈里，斯蒂芬A。(1998)。天文学与天体物理学导论(第四版)。桑德斯学院出版。第56页。ISBN 978-0-03-006228-5。
\item 威廉姆斯，大卫R.(2006年2月10日)。"行星概况"。NASA.请参阅太阳和月亮页面上的表观直径。已存档从2016年3月4日的原件。已检索9月28日2008。
\item 威廉姆斯，大卫R.(2004年9月1日)。"月球概况介绍"。NASA.已存档从原件到2020年6月13日。已检索3月21日2007。
\item 巴斯克斯，M.；罗德里格斯，P。蒙特涅斯; 帕勒，E。(2006)。“在寻找太阳系外行星时，地球是天体物理学感兴趣的对象”(PDF)。天体物理学讲稿和论文。2: 49。Bibcode:2006LNEA....2...49V。存档自原来的(PDF)2011年8月17日。已检索3月21日2007。
\item 天体物理学家团队 (2005年12月1日)。“地球在银河系中的位置”。NASA.存档自原来的2008年7月1日。已检索6月11日2008。
\item Rohli, Robert.五、；维加，安东尼J.(2018)。气候学(第四版)。琼斯和巴特利特学习。pp. 291-292。ISBN 978-1-284-12656-3。
\item 克里斯·伯恩 (1996年3月)。极地之夜(PDF)。奥罗拉研究所。已存档(PDF)从2023年8月6日开始。已检索9月28日2015。
\item "日照时间"。澳大利亚南极计划。2020年6月24日。已存档从原件到2020年10月22日。已检索10月13日2020。
\item Bromberg，Irv (2008年5月1日)。“四季的长度 (在地球上)”。Sym545。多伦多大学。存档自原来的2008年12月18日。已检索11月8日2008。
\item 林浩生 (2006)。“月球轨道进动的动画”。天文学AST110-6调查。夏威夷大学马诺阿分校。已存档从原来的2010年12月31日。已检索9月10日2010。
\item 里克·费舍尔 (1996年2月5日)。“地球自转和赤道坐标”。国家射电天文台。存档自原来的2011年8月18日。已检索3月21日2007。
\item Buis，Alan (2020年2月27日)。“米兰科维奇 (轨道) 周期及其在地球气候中的作用”。NASA。已存档从原件到2020年10月30日。已检索10月27日2020。
\item 莎拉·M·康；西格，理查德。“Croll再访: 为什么北半球比南半球温暖？”(PDF)。哥伦比亚大学。纽约。已存档(PDF)从原始的2021年9月7日。已检索10月27日2020。
\item 克莱梅蒂，埃里克 (2019年6月17日)。“我们的月亮到底有什么特别之处？”。天文学。已存档从原件到2020年11月6日。已检索10月13日2020。
\item "Charon"。NASA。2019年12月19日。已存档从原件到2020年10月14日。已检索10月13日2020。
\item 布朗，托比 (2019年12月2日)。好奇的孩子: 为什么月亮叫月亮？。对话。已存档从原件到2020年11月8日。已检索10月13日2020。
\item Chang，Kenneth (2023年11月1日)。一项研究表明，一个 “大怪兽” 形成了月球，并在地球深处留下了痕迹-地球深处的两个巨大斑点可能是月球诞生的残余物。。纽约时报。已存档从2023年11月1日开始。已检索11月2日2023。
\item 袁，钱; 等 (2023年11月1日)。“形成月球的撞击器是地球基底地幔异常的来源”。自然。623(7985):95-99。Bibcode:2023年自然623. .. 95Y。doi:10.1038/s41586-023-06589-1。PMID37914947。S2CID264869152。已存档从2023年11月2日的原件。已检索11月2日2023。  
\item 咳嗽，克里斯托弗·L；阿彻，艾伦·W；肯尼斯·J·拉科瓦拉.(2009)。“地球-月球系统中的潮汐，潮汐和长期变化”。地球科学评论。97(1):59-79.Bibcode:2009ESRv...97...59C。doi:10.1016/j.earscirev.2009.09.002。ISSN0012-8252。已存档从2012年10月28日的原件。已检索10月8日2020。 
\item 凯利，彼得 (2017年8月17日)。“潮汐锁定的系外行星可能比以前认为的更常见”。Uw新闻。已存档从原件到2020年10月9日。已检索10月8日2020。
\item “月相和月食 | 地球的月亮”。NASA太阳系探索。已存档从原件到2020年10月16日。已检索10月8日2020。
\item 弗雷德·埃斯佩纳克;吉恩·米乌斯(2007年2月7日)。“月亮的长期加速”。NASA.存档自原来的2008年3月2日。已检索4月20日2007。
\item 威廉姆斯，g.e.(2000)。“地球自转和月球轨道的前寒武纪历史的地质限制”。地球物理学的评论。38(1):37-59。Bibcode:2000RvGeo .. 38...37W。doi:10.1029/1999RG900016。S2CID51948507。 
\item Laskar，J.; 等人 (2004年)。“地球日照量的长期数值解”。天文学和天体物理学。428(1):261-285。Bibcode:2004A & A...428 .. 261L。doi:10.1051/0004-6361:20041335。已存档从原件到2018年5月17日。已检索5月16日2018。
\item 库珀，基思 (2015年1月27日)。“地球的月亮可能对生命并不重要”。Phys.org。已存档从原件到2020年10月30日。已检索10月26日2020。
\item 达达里奇，艾米；米特罗维察，杰里·X。松山，Isamu; 佩伦，J。泰勒；迈克尔·漫画理查兹，马克·A。(2007年11月22日)。“平衡旋转稳定性与火星图形”(PDF)。伊卡洛斯。194(2):463-475。doi:10.1016/j.icarus.2007.10.017。存档自原来的(PDF)2020年12月1日。已检索10月26日2020。
\item Sharf, Caleb A.(2012年5月18日)。“日食巧合”。科学美国人。已存档从原件到2020年10月15日。已检索10月13日2020。
\item 克里斯托·阿波斯托洛斯·A；亚瑟，大卫·J.(2011年3月31日)。“地球的长寿马蹄形伴侣”。皇家天文学会月刊。414(4):2965-2969。arXiv:1104.0036。Bibcode:2011MNRAS.414.2965C。doi:10.1111/j.1365-2966.2011.18595.x。S2CID13832179。 见表2，第5页。
\item Marcos，C. de la Fuente; Marcos，R. de la Fuente (2016年8月8日)。“小行星 (469219) 2016 HO3，最小，最接近的地球准卫星”。皇家天文学会月刊。462(4):3441-3456。arXiv:1608.01518。Bibcode:2016MNRAS.462.3441D。doi:10.1093/mnras/stw1972。S2CID118580771。已存档从原件到2020年10月31日。已检索10月28日2020。 
\item Choi, Charles Q.(2011年7月27日)。“终于发现了地球的第一颗小行星伴侣”。Space.com。已存档从原来的2013年8月26日。已检索7月27日2011年。
\item “2006 RH120 ( = 6 r10db9) (地球的第二个月亮？)”。大谢福德天文台。存档自原来的2015年2月6日。已检索7月17日2015。
\item "UCS卫星数据库"。核武器与全球安全。关注科学家联盟。2021年9月1日。已存档从原来的2016年1月25日。已检索1月12日2022年。
\item 韦尔奇，罗珊娜; 拉姆菲尔，佩格A。(2019年)。美国历史上的技术创新: 科学技术百科全书 [3卷]。ABC-克里奥。第126页。ISBN 978-1-61069-094-2。已存档从2023年8月10日开始。已检索8月9日2023。
\item 马修·夏雷特；史密斯，沃尔特H. F.(2010年6月)。“地球海洋的体积”。海洋学。23(2):112-114。Bibcode:2010Ocgpy .. 23 b.112C。doi:10.5670/oceanog.2010.51。高密度脂蛋白:1912/3862。
\item “来自太阳的第三块岩石 -- 不安分的地球”。NASA的宇宙。已存档从2015年11月6日的原件。已检索4月12日2015。
\item 欧洲投资银行 (2019)。在水上。出版物办公室。doi:10.2867/509830。ISBN 978-92-861-4319-9。已存档从原件到2020年11月29日。已检索12月7日2020。
\item Khokhar，Tariq (2017年3月22日)。图表: 全球70\% 的淡水用于农业。世界银行博客。已存档从原件到2020年12月6日。已检索12月7日2020。
\item Perlman，Howard (2014年3月17日)。“世界之水”。USGS水科学学校。已存档从2015年4月22日的原件。已检索4月12日2015。
\item “湖在哪里？”。湖泊科学家。2016年2月28日已存档从原始的2023年2月28日。已检索2月28日2023。
\item  学校，水科学 (2019年11月13日)。“地球上有多少水？-美国地质调查局”。USGS.gov。已存档从2022年6月9日开始。已检索3月3日2023。
\item “淡水资源”。教育。2022年8月18日。已存档从2022年5月26日开始。已检索2月28日2023。
\item 亨德里克斯，马克 (2019)。地球科学: 导论。波士顿: Cengage。第330页。ISBN 978-0-357-11656-2。
\item 亨德里克斯，马克 (2019)。地球科学: 导论。波士顿: Cengage。第329页。ISBN 978-0-357-11656-2。
\item 肯尼什，迈克尔·J。(2001)。海洋科学实用手册。海洋科学丛书 (第三版)。佛罗里达州博卡拉顿: CRC出版社。第35页。doi:10.1201/9781420038484。ISBN 978-0-8493-2391-1。
\item Mullen，Leslie (2002年6月11日)。“早期地球的盐”。NASA天体生物学杂志。存档自原来的2007年6月30日。已检索3月14日2007。
\item 罗恩·M·莫里斯.“海洋过程”。NASA天体生物学杂志。存档自原来的2009年4月15日。已检索3月14日2007。
\item Michon Scott (2006年4月24日)。“地球的大热水桶”。NASA地球观测站。已存档从2008年9月16日开始。已检索3月14日2007。
\item 样本，Sharron (2005年6月21日)。“海面温度”。NASA.存档自原来的2013年4月27日。已检索4月21日2007。
\item 中心，天体地质科学 (2021年10月14日)。“太阳系中的水之旅-美国地质调查局”。USGS.gov。已存档从2022年1月19日开始。已检索1月19日2022年。
\item “其他星球上有海洋吗？”。NOAA国家海洋服务。2013年6月1日。已存档从原件到2017年6月19日。已检索1月19日2022年。
\item Exline，约瑟夫·D；莱文，阿琳·S；莱文，乔尔·S.(2006)。气象学: 5-9年级基于探究的学习的教育者资源(PDF)。NASA/兰利研究中心。第6页。NP-2006-08-97-LaRC。已存档(PDF)从原件到2018年5月28日。已检索7月28日2018。
\item 金，迈克尔·D；普拉特尼克，史蒂文; 门泽尔，W。保罗; 阿克曼，史蒂文A；胡班克斯，保罗·A。(2013)。“Terra和Aqua卫星上的MODIS观测到的云的时空分布”。IEEE地球科学与遥感汇刊。51(7)。电气和电子工程师协会 (IEEE):3826-3852。Bibcode:2013ITGRS..51.3826K。doi:10.1109/tgrs.2012.2227333。高密度脂蛋白:2060/20120010368。ISSN0196-2892。S2CID206691291。  
\item Geerts, B.;Linacre, E.(1997年11月)。“对流层顶的高度”。大气科学资源。怀俄明大学。已存档从原件到2020年4月27日。已检索8月10日2006。
\item 哈里森，罗伊M.；罗纳德·E·海丝特。(2002)。Uv-b辐射增加的原因和环境影响。皇家化学学会。ISBN 978-0-85404-265-4。
\item 工作人员 (2003年10月8日)。“地球的大气层”。NASA.存档自原来的2020年4月27日。已检索3月21日2007。
\item 迈克尔·皮德维尼 (2006)。《自然地理学基础 (第2版) 》。不列颠哥伦比亚大学，奥肯那根。已存档从2011年8月15日的原件。已检索3月19日2007。
\item Gaan，Narottam (2008)。气候变化与国际政治。Kalpaz出版物。第40页。ISBN 978-81-7835-641-9。
\item 纳迪亚·德雷克(2018年12月20日)。“空间的边缘到底在哪里？这取决于你问谁”。国家地理。存档自原来的2021年3月4日。已检索12月4日2021。
\item 埃里克森，克里斯汀; 道尔，希瑟 (2019年6月28日)。“对流层”。SpacePlace。NASA。已存档从原件到2021年12月4日。已检索12月4日2021。
\item 莫兰，约瑟夫·M。(2005)。"天气"。世界图书在线参考中心。NASA/世界图书公司存档原来的2010年12月13日。已检索3月17日2007。
\item 伯杰，沃尔夫冈H.(2002)。“地球的气候系统”。加州大学圣地亚哥分校。已存档从2013年3月16日的原件。已检索3月24日2007。
\item 斯特凡·拉姆斯托夫(2003)。“温盐海洋环流”。波茨坦气候影响研究所。已存档从2013年3月27日的原件。已检索4月21日2007。
\item 科丁顿，奥德勒；精益，朱迪思·L。皮尔斯基，彼得; 斯诺，马丁; 林德霍尔姆，道格 (2016)。“太阳辐照度气候数据记录”。美国气象学会公报。97(7):1265-1282。Bibcode:2016BAMS...97.1265C。doi:10.1175/bams-d-14-00265.1。
\item 大卫·萨达瓦；海勒，H.克雷格；Orians, Gordon H.(2006)。生命，生物学的科学(第8版)。麦克米伦.p. 1114。ISBN 978-0-7167-7671-0。
\item 工作人员。“气候区”。英国环境、食品和农村事务部。存档自原来的2010年8月8日。已检索3月24日2007。
\item Rohli, Robert.五、；维加，安东尼J.(2018)。气候学(第四版)。琼斯和巴特利特学习。第49页。ISBN 978-1-284-12656-3。
\item Rohli, Robert.五、；维加，安东尼J.(2018)。气候学(第四版)。琼斯和巴特利特学习。第32页。ISBN 978-1-284-12656-3。
\item Rohli, Robert.五、；维加，安东尼J.(2018)。气候学(第四版)。琼斯和巴特利特学习。第34页。ISBN 978-1-284-12656-3。
\item Rohli, Robert.五、；维加，安东尼J.(2018)。气候学(第四版)。琼斯和巴特利特学习。第46页。ISBN 978-1-284-12656-3。
\item 各种 (1997年7月21日)。“水文循环”。伊利诺伊大学。已存档从2013年4月2日的原件。已检索3月24日2007。
\item Rohli, Robert.五、；维加，安东尼J.(2018)。气候学(第四版)。琼斯和巴特利特学习。第159页。ISBN 978-1-284-12656-3。
\item El Fadli，Khalid I.; 等人 (2013年)。“世界气象组织对利比亚El Azizia所谓的世界纪录58 °C极端温度的评估 (1922年9月13日)”。美国气象学会公报。94(2):199-204.Bibcode:2013BAMS...94..199E。doi:10.1175/BAMS-D-12-00093.1。ISSN0003-0007。 
\item 特纳，约翰; 等人 (2009)。“创纪录的南极沃斯托克站地面气温低”。地球物理研究杂志: 大气。114(D24) 2009JD012104: d24102。Bibcode:2009JGRD..11424102T。doi:10.1029/2009JD012104。ISSN2156-2202。 
\item 莫顿，奥利弗 (2022年8月26日)。"上层大气定义"。柯林斯·w ö rterbuch(德语)。已存档从2023年2月21日开始。已检索8月26日2022年。
\item 工作人员 (2004年)。“平流层和天气; 平流层的发现”。科学周。存档自原来的2007年7月13日。已检索3月14日2007。
\item 德科尔多瓦，S.桑兹·费尔南德斯 (2004年6月21日)。“卡曼分隔线的介绍，用作分隔航空和航天的边界”。国际航空联合会。存档自原来的2010年1月15日。已检索4月21日2007。
\item 刘，S。C.;多纳休，T。M。(1974)。“地球大气中氢的空气动力学”。大气科学杂志。31(4):1118-1136。Bibcode:1974JAtS。31.1118l。doi:10.1175/1520-0469(1974)031<1118:TAOHIT>2.0.CO;2。
\item 卡特林，大卫·C.;扎恩勒，凯文J。;麦凯，克里斯托弗·P。(2001)。“生物甲烷、氢逃逸和早期地球的不可逆氧化”。科学。293(5531):839-843。Bibcode:2001Sci...293 .. 839C。CiteSeerX10.1.1.562.2763。doi:10.1126/science.1061976。PMID11486082。S2CID37386726。   
\item 阿贝登，斯蒂芬·T.(1997年3月31日)。“地球的历史”。俄亥俄州立大学。存档自原来的2012年11月29日。已检索3月19日2007。
\item 亨顿，D。M、；多纳休，T。M(1976)。“地球行星的氢损失”。地球与行星科学年度评论。4(1):265-292。Bibcode:1976AREPS...4..265H。doi:10.1146/annurev.ea.04.050176.001405。
\item 工作人员 (2003年9月)。“天体生物学路线图”。美国航空航天局，洛克希德·马丁公司。存档自原来的2012年3月12日。已检索3月10日2007。
\item 辛格，J.S。;辛格，S.P。；古普塔，S.R.(2013)。生态环境科学与保护(第一版)。新德里: S.钱德公司。ISBN 978-93-83746-00-2。OCLC 896866658。
\item Rutledge，Kim; 等人 (2011年6月24日)。“生物圈”。国家地理。已存档从原始的2022年5月28日。已检索11月1日2020。
\item “动物和植物物种之间的相互依存关系”。BBC Bitesize。英国广播公司。P.3。已存档从原件到2019年6月27日。已检索6月28日2019。
\item 赫尔穆特·希勒布兰德 (2004年)。“关于纬度梯度的一般性”(PDF)。美国博物学家。163(2):192-211。doi:10.1086/381004。PMID14970922。S2CID9886026。已存档(PDF)从2017年9月22日的原件。已检索4月20日2018。   
\item “NASA天体生物学研究所”。astrobiology.nasa.gov。已存档从原件到2023年11月17日。已检索11月9日2023。
\item 沙龙·史密斯弗莱明，劳拉; 索洛-加布里埃尔，海伦娜; 格里克，威廉·H (2011)。海洋与人类健康。爱思唯尔科学。第212页。ISBN 978-0-08-087782-2。
\item 亚历山大，大卫 (1993)。自然灾害。施普林格科学与商业媒体。第3页。ISBN 978-1-317-93881-1。已存档从2023年8月10日开始。已检索8月9日2023。
\item “什么是气候变化？”。联合国。已存档从2023年1月26日开始。已检索8月17日2022年。
\item 安德鲁·古迪(2000)。人类对自然环境的影响。麻省理工学院出版社。第52、66、69、137、142、185、202、355、366页。ISBN 978-0-262-57138-8。
\item 厨师，约翰；Oreskes, Naomi;多兰，彼得T.；威廉·R·安德雷格。L.;巴特·维赫根；艾德·W·梅巴赫；卡尔顿，J。斯图尔特；斯蒂芬·莱万多夫斯基；安德鲁·G·斯库斯；格林，莎拉A；Nuccitelli，Dana; Jacobs，Peter; Richardson，Mark; Winkler，b ä rbel; 绘画，Rob; Rice，Ken (2016)。“共识共识: 对人类造成的全球变暖的共识估计的综合”。环境研究信函。11(4) 048002。Bibcode:2016ERL .... 11d8002C。doi:10.1088/1748-9326/11/4/048002。高密度脂蛋白:1983/34949783-dac1-4ce7-ad95-5dc0798930a6。ISSN1748-9326。 
\item “全球变暖效应”。国家地理。2019年1月14日。存档自原来的2017年1月18日。已检索9月16日2020。
\item “人类进化导论 | 史密森学会人类起源计划”。humanorigins.si.edu。2022年7月11日。已存档从2023年11月8日开始。已检索11月9日2023。
\item “80亿日”。联合国。2025年11月15日已存档从2025年5月23日开始。已检索6月23日2025。
\item 菲奥娜·哈维(2020年7月15日)。研究表明，2100年世界人口可能比联合国预测低20亿。守护者。ISSN0261-3077。已存档从原件到2020年9月4日。已检索9月18日2020。 
\item Lutz，Ashley (2012年5月4日)。每日地图: 几乎每个人都生活在北半球。商业内幕。已存档从原件到2018年1月19日。已检索1月5日2019。
\item 阿贝尔·门德斯(2011年7月6日)。“古地球陆地的分布”。波多黎各大学阿雷西博分校。存档自原来的2019年1月6日。已检索1月5日2019。
\item 里奇，H。;罗瑟，M.(2019年)。“未来将有多少人生活在城市地区？”。我们的数据世界。已存档从原件到2020年10月29日。已检索10月26日2020。
\item 大卫·谢勒; 伯特·维斯 (2005年)。俄罗斯宇航员: 在尤里·加加林训练中心内。伯克哈瑟.ISBN 978-0-387-21894-6。
\item 福尔摩斯，奥利弗 (2018年11月19日)。“太空: 我们走了多远-我们要去哪里？”。守护者。ISSN0261-3077。已存档从原件到2020年10月6日。已检索10月10日2020。 
\item “会员国 | 联合国”。联合国。已存档从2023年3月1日开始。已检索1月3日2024。
\item 约翰·劳埃德;约翰·米奇森(2010)。离散丰美的第二气一般无知之书。Faber & Faber. pp. 116-117。ISBN 978-0-571-29072-7。
\item 史密斯，考特尼B。(2006)。联合国的政治与进程: 全球之舞(PDF)。Lynne Reiner. pp. 1-4.ISBN  978-1-58826-323-0。已存档 (PDF)从原件到2020年10月17日。已检索10月14日2020。
\item “过度开发自然资源的后果是什么？”。Iberdrola。已存档从原件到2019年6月27日。已检索6月28日2019。
\item ”13。“开发自然资源”。欧洲环境署。欧洲联盟。2016年4月20日已存档从原件到2019年6月27日。已检索6月28日2019。
\item Huebsch，罗素 (2017年9月29日)。“如何从地下提取化石燃料？”。科学。“叶” 组媒体。已存档从原件到2019年6月27日。已检索6月28日2019。
\item “发电-有什么选择？”。世界核协会。已存档从原件到2019年6月27日。已检索6月28日2019。
\item 乔治·布里姆霍尔 (1991年5月)。“矿石的成因”。科学美国人。264(5)。自然美国:84-91。Bibcode:1991 SciAm.264e .. 84B。doi:10.1038/scientificamerican0591-84。JSTOR24936905。 
\item 鲁宁，乔纳森一世。(2013)。地球: 可居住世界的演变(第二版)。剑桥大学出版社。pp. 292-294.ISBN 978-0-521-61519-8。
\item 罗娜，彼得A.(2003)。“海底资源”。科学。299(5607):673-674。doi:10.1126/science.1080679。PMID12560541。S2CID129262186。  
\item 里奇，H。;罗瑟，M.(2019年)。“土地使用”。我们的数据世界。已存档从2019年4月11日的原件。已检索10月26日2020。
\item 气专委(2019年)。“决策者摘要”(PDF)。IPCC关于气候变化和土地的特别报告。第8页。已存档(PDF)从原件到2020年2月17日。已检索9月25日2020。
\item 尼基·泰特塔特-斯特拉顿，丹妮 (2014)。避难: 在世界各地的家中。虎鲸图书出版商。第6页。ISBN 978-1-4598-0742-6。
\item 气专委(2021)。马森-德尔莫特，V.；翟，P.；皮拉尼，A.；康纳斯，S.等。(编辑)。2021年气候变化: 物理科学基础(PDF)。第一工作组对第六次评估报告政府间气候变化专门委员会。英国剑桥和美国纽约: 剑桥大学出版社 (出版中)。SPM-7。已存档(PDF)从原件到2021年8月13日。已检索6月2日2022年。
\item Lindsey，Rebecca (2009年1月14日)。“气候和地球的能源预算”。地球观测站。NASA。已存档从2019年10月2日的原件。已检索12月19日2021。
\item 《2020年全球气候状况》。世界气象组织。2021年1月14日存档自原来的2023年11月29日。已检索3月3日2021。
\item DiGirolamo，迈克 (2021年9月8日)。“我们已经跨越了九个行星边界中的四个。这是什么意思？”。Mongabay。已存档从最初的2022年1月27日。已检索1月27日2022年。
\item 卡灵顿，达米恩 (2022年1月18日)。“化学污染已经超过了人类的安全极限，科学家说”。守护者。已存档从2022年4月12日开始。已检索1月27日2022年。
\item 丹尼尔·W·奥尼尔；范宁，安德鲁·L；羔羊，威廉·F；斯坦伯格，朱莉娅K.(2018)。“地球边界内所有人的美好生活”。自然可持续性。1(2):88-95。Bibcode:2018NatSu...1...88O。doi:10.1038/s41893-018-0021-4。ISSN2398-9629。S2CID169679920。已存档从2022年2月1日开始。已检索1月30日2022年。  
\item 泰德·威德默(2018年12月24日)。柏拉图认为地球是什么样子？几千年来，人类一直试图想象太空中的世界。五十年前，我们终于看到了它。。纽约时报。存档自原来的2022年1月1日。已检索12月25日2018。
\item 卡尔·G·利翁格曼.(2004)。“第29组: 多轴对称，软线和直线，带交叉线的封闭标志”。符号-西方符号和表意符号百科全书。纽约: Ionfox AB. pp。 281-282。ISBN 978-91-972705-0-2。
\item Stokey，Lorena Laura (2004)。世界神话专题指南。韦斯特波特，CN: 格林伍德出版社。pp. 114-115。ISBN 978-0-313-31505-3。
\item 詹姆斯·E·洛夫洛克.(2009)。盖亚消失的脸。基本书籍。第255页。ISBN 978-0-465-01549-8。
\item 詹姆斯·E·洛夫洛克.(1972)。“透过大气层看到的盖亚”。大气环境。6(8):579-580。Bibcode:1972阿特曼。6。579L。doi:10.1016/0004-6981(72)90076-5。ISSN1352-2310。 
\item 洛夫洛克，J.E.；Margulis, L.(1974)。“生物圈的大气稳态: 盖亚假说”。特勒斯A。26(1-2):2-10。Bibcode:1974告诉... 26 .... 2L。doi:10.3402/tellusa.v26i1-2.9731。S2CID129803613。 
\item 再见，丹尼斯(2018年12月21日)。“阿波罗8号的地球崛起: 环游世界的镜头-半个世纪前的今天，一张来自月球的照片帮助人类重新发现了地球”。纽约时报。存档自原来的2022年1月1日。已检索12月24日2018。
\item 博尔顿，马修·迈尔; 海特豪斯，约瑟夫 (2018年12月24日)。“我们都是同一个星球上的骑手 -- 50年前从太空看，地球是一份值得保存和珍惜的礼物。怎么了？”。纽约时报。存档自原来的2022年1月1日。已检索12月25日2018。
\item “ESPI晚间活动” 看到我们的整个星球: 地球观测的文化和伦理观点"。ESPI-欧洲空间政策研究所。2021年10月7日。已存档从最初的2022年1月27日。已检索1月27日2022年。
\item “两张支持环保运动的地球早期图像-CBC电台”。CBC。2020年4月16日。已存档从最初的2022年1月27日。已检索1月27日2022年。
\item 布雷，凯伦; 伊顿，希瑟; 鲍曼，惠特尼，编辑。(2023年10月3日)。尘世的事物: 内在性，新唯物主义和行星思维。福特汉姆大学出版社。doi:10.5422/福特汉姆/9781531503055.001.0001。ISBN 978-1-5315-0413-7。
\item 卡恩，查尔斯H。(2001)。毕达哥拉斯和毕达哥拉斯: 简史。印第安纳波利斯，在剑桥，英国: 哈克特出版公司。第53页。ISBN 978-0-87220-575-8。已存档从原件到2023年12月14日。已检索8月9日2023。
\item 克里斯汀·加伍德 (2008)。平坦的地球: 一个臭名昭著的想法的历史(第一版)。纽约: 托马斯·邓恩图书。pp。 26-31.ISBN 978-0-312-38208-7。OCLC 184822945。
\item “地心模型 | 定义、历史和事实 | 大英百科全书”。大英百科全书。已存档从2025年9月9日开始。已检索12月2日2025。
\item Arnett，Bill (2006年7月16日)。"地球"。九大行星，太阳系的多媒体之旅: 一颗恒星，八颗行星等等。已存档从2000年8月23日起。已检索3月9日2010。
\item 詹姆斯·门罗; 里德·威坎德; 理查德·黑兹利特 (2007)。物理地质学: 探索地球。汤姆森·布鲁克斯/科尔。pp. 263-265。ISBN 978-0-495-01148-4。
\item 亨肖，约翰·M。(2014)。每个场合的方程式: 52个公式和它们为什么重要。约翰霍普金斯大学出版社.pp. 117-118.ISBN 978-1-4214-1491-1。
\item 乔·D·伯奇菲尔德.(1990)。开尔文勋爵和地球的年龄。芝加哥大学出版社。pp. 13-18.ISBN 978-0-226-08043-7。
\end{enumerate}
\subsection{外部链接}
\begin{itemize}
\item 大地剖面-太阳系探索-NASA
\item 地球观测站-NASA
\item 地球-视频-国际空间站:
\item 视频 (01:02)在YouTube上-地球 (延时)
\item Video (00:27) on YouTube – Earth and auroras (time-lapse)
\item Google Earth 3D, interactive map
\item Interactive 3D visualization of the Sun, Earth and Moon system
\item GPlates Portal (University of Sydney)
\end{itemize}